\documentclass[
  reprint,
  amsmath,
  amssymb,
  aps,
  pra,
  nofootinbib,
  superscriptaddress,
  floatfix
]{revtex4-2}

\usepackage[T1]{fontenc}
\usepackage{lmodern}
\usepackage{graphicx}
\usepackage{capt-of}
\graphicspath{{../../}{}}
\usepackage{booktabs}
\usepackage{bm}
\usepackage{mathtools}
\usepackage{xcolor}
\usepackage{hyperref}

\newcommand{\placeholder}[1]{\textcolor{red!70!black}{\ifmmode\left[#1\right]\else\textsf{[#1]}\fi}}
\newcommand{\tentative}[1]{\textcolor{red!70!black}{#1}}
\newcommand{\tdtrip}[3]{\shortstack[c]{#1/#2\\/#3}}
\newcommand{\tdacc}[2]{\shortstack[c]{(#1,\\#2)}}
\newcommand{\tdsuz}[3]{\shortstack[c]{(#1,#2);\\$D=#3$}}
\newcommand{\metrichead}[1]{\makebox[0.052\textwidth][c]{#1}}
\newcommand{\trajectoryplaceholder}[1]{%
  \begin{center}
  \fbox{\begin{minipage}[c][1.35in][c]{0.94\columnwidth}
  \centering\small\tentative{#1 trajectory panel}\\[0.6ex]
  \scriptsize ED reference, AP-McLachlan, pinned Qiskit-community TrotterQRTE, PVQD, VarQRTE
  \end{minipage}}
  \end{center}
}
\newcommand{\structlabel}[1]{\noindent\mbox{\tentative{#1.}}~}
\newcommand{\Oset}{\mathcal{O}}
\newcommand{\Gset}{\mathcal{G}}
\newcommand{\Rset}{\mathcal{R}}
\newcommand{\Sset}{\mathcal{S}}
\newcommand{\Hcal}{\mathcal{H}}
\newcommand{\Bcal}{\mathcal{B}}
\newcommand{\Ocal}{\mathcal{O}}
\newcommand{\Ucal}{\mathcal{U}}
\newcommand{\Jcal}{\mathcal{J}}
\newcommand{\Pcal}{\mathcal{P}}
\newcommand{\Dcal}{\mathcal{D}}
\newcommand{\Mcal}{\mathcal{M}}
\newcommand{\Wcal}{\mathcal{W}}
\newcommand{\dd}{\mathrm{d}}
\newcommand{\ii}{\mathrm{i}}
\newcommand{\ee}{\mathrm{e}}
\newcommand{\TR}{\mathrm{TR}}
\newcommand{\ex}{\mathrm{ex}}
\newcommand{\ctrl}{\mathrm{ctrl}}
\newcommand{\stat}{\mathrm{stat}}
\newcommand{\dyn}{\mathrm{dyn}}
\newcommand{\drv}{\mathrm{drv}}
\newcommand{\obs}{\mathrm{obs}}
\newcommand{\stay}{\mathrm{stay}}
\newcommand{\app}{\mathrm{append}}
\newcommand{\prune}{\mathrm{prune}}
\newcommand{\swap}{\mathrm{exchange}}
\newcommand{\Var}{\operatorname{Var}}
\newcommand{\Tr}{\operatorname{Tr}}
\newcommand{\Top}{\operatorname{Top}}
\newcommand{\argminp}{\mathop{\mathrm{arg\,min}}}
\newcommand{\argmaxp}{\mathop{\mathrm{arg\,max}}}
\newcommand{\proj}{\mathrm{proj}}
\newcommand{\diag}{\mathrm{diag}}
\newcommand{\CVaR}{\operatorname{CVaR}}
\newcommand{\softmax}{\operatorname{softmax}}
\newcommand{\hh}{\mathrm{HH}}
\newcommand{\phmax}{n_{\mathrm{ph,max}}}
\newcommand{\twq}{\mathrm{2Q}}
\newcommand{\eps}{\varepsilon}
\newcommand{\ket}[1]{|#1\rangle}
\newcommand{\bra}[1]{\langle #1|}
\newcommand{\braket}[2]{\langle #1|#2\rangle}
\newcommand{\expval}[1]{\langle #1\rangle}
\newcommand{\inlinetablecaption}[2]{%
  \refstepcounter{table}\label{#1}%
  \noindent\textsc{Table~\thetable.} #2\par\vspace{0.6ex}%
}


\usepackage{multirow}
\begin{document}

\title{Append--Prune McLachlan Dynamics}

\author{Jake Strobel}
\author{Jacopo Simoni}
\author{Yuan Ping}
\affiliation{University of Wisconsin--Madison}

\date{\today}


\begin{abstract}
Projected variational dynamics can fail because the active tangent manifold does not represent the driven Schr\"odinger drift or because the numerical solve and finite time step do not realize motion that the manifold already contains.  We present Append--Prune McLachlan dynamics (AP-McLachlan), a checkpoint-local policy that diagnoses these structural and numerical failures separately.  Its structural decision is one generalized exchange patch with removal and addition components.  Empty components give stay, pure append, or pure prune as one-sided limits; two nonempty components give true delete--append exchange.  Candidate patches are enumerated as complete deletion-cardinality rungs paired with positioned zero-angle insertions at commutation-reduced circuit cuts, scored by realized captured drift per compiled hardware cost, and certified one finalist at a time on the complete patched tangent support through conditioning, ray-continuity, and velocity-smoothness gates with a bounded trust-region refit.  When the McLachlan distance exceeds its declared cut, new tangent measurements open the full generalized-exchange family; measurement-free pruning remains available at every checkpoint.  On the accepted fixed support, a state-space guard subdivides the integration step whenever the projected state displacement over the step exceeds a declared bound; a wider inverse-candidate search is retained as a diagnostic for pathological checkpoints but is not the operating path.  Exact trajectories remain post-run references and never enter support selection, solve repair, or time-step decisions.  Driven Hubbard--Holstein diagnostics isolate the numerical-repair, append, and prune mechanisms through same-seed energy, doublon, site-occupation, residual, and compiled-resource comparisons; quantitative cross-regime claims remain \tentative{to be transferred from the source-locked benchmark record}.
\end{abstract}

\maketitle

\section{Introduction}

\subsection{Modeled problem}
Driven mixed fermion--boson systems require simultaneous propagation of electronic correlation, bosonic response, and their mutual dressing \cite{Macridin2018,Macridin2018PRA,Sawaya2020,Li2023,Denner2023}.  The physical outputs are trajectories: energy absorption, doublon response, total and spin-resolved site occupations, density modulation, and frequency-domain response.  We use the driven Hubbard--Holstein Hamiltonian
\begin{align}
H_{\rm HH}(t)
&=-t\sum_{\langle i,j\rangle,\sigma}
(c_{i\sigma}^{\dagger}c_{j\sigma}+c_{j\sigma}^{\dagger}c_{i\sigma})
+U\sum_i n_{i\uparrow}n_{i\downarrow}
\nonumber\\
&\quad+\omega_0\sum_i b_i^\dagger b_i
+g\sum_i(b_i^\dagger+b_i)(n_i-\bar n)
\nonumber\\
&\quad+\sum_{i,\sigma}v_i(t)n_{i\sigma},
\nonumber\\
v_i(t)
&=s_i A\sin(\omega t+\phi)
\nonumber\\
&\quad\times
\exp\left[-\frac{(t-t_0)^2}{2\bar t^{\,2}}\right].
\label{eq:hubbard_holstein_driven}
\end{align}
Here \(c_{i\sigma}\) and \(b_i\) correspond to fermionic and bosonic annihilation operators respectively, acting on site $i$, with creation given by $\dagger$; $\sigma \in\{\uparrow, \downarrow\}$ denotes particle spin;
\(n_{i\sigma}, n_i\) are the spin-resolved and total fermionic number operators respectively; \(\bar n=N_e/L\) is the mean fermion filling
for fixed electron number \(N_e\) and lattice size \(L\), and the site profile \(s_i\) selects the driven density mode. Our results use open boundary conditions with half filling.  The near-term objective is to track these observables without allowing circuit depth, estimator work, or tangent-metric conditioning to destroy the propagated trajectory.

\subsection{Current status: projective McLachlan dynamics}
Product formulas approximate the full unitary evolution directly and provide a depth--accuracy baseline \cite{Suzuki1985,Suzuki1990,ChildsSu2019ProductFormula,ChildsOstranderSu2019RandomizedPF,Campbell2019QDRIFT}.  Variational alternatives propagate a parameterized state on a restricted circuit manifold.  Projective McLachlan simulation, pVQD, adaptive pVQD, AVQDS, and compressed adaptive dynamics establish fixed-manifold propagation and residual- or projection-driven circuit growth \cite{McLachlan1964,LiBenjamin2017VQS,Yuan2019VQS,Endo2020GeneralProcesses,Benedetti2021TimeEvolution,Barison2021,Yao2021,Linteau2024,Zhang2025}.  These methods establish adaptive growth and zero-initialized additions as prior art.  Mixed-system spin--boson benchmarks further show that compiled cost must be interpreted together with trajectory accuracy and noise sensitivity \cite{Miessen2021SpinBoson}.

At time \(t_k\), standard projective McLachlan propagation chooses the tangent velocity that best approximates the horizontally projected Schr\"odinger drift.  If \(\bar T_k\) is the projected ansatz tangent block and \(\bar b_k\) is the projected physical drift, then
\begin{equation}
\dot\theta_k
\in\operatorname*{arg\,min}_{v}
\|\bar T_kv-\bar b_k\|_2^2,
\qquad
G_k\dot\theta_k=f_k,
\label{eq:mclachlan_velocity}
\end{equation}
where \(G_k=\Re(\bar T_k^\dagger\bar T_k)\) and \(f_k=\Re(\bar T_k^\dagger\bar b_k)\).  Equation~(\ref{eq:mclachlan_velocity}) defines the standard propagation problem.  AP-McLachlan changes which tangent directions are active and how the resulting finite-precision solve is realized; it does not change the McLachlan principle.

\subsection{Limitations with current approaches}
The remaining failure has two components.  First, the manifold can be structurally wrong: useful tangent directions may be absent, while obsolete or mutually redundant directions remain.  A compact ADAPT ansatz may prepare the initial state accurately yet fail to span the later driven velocity \cite{Grimsley2019,Tang2021,Yordanov2021}.  Append-only dynamics addresses missing directions but permits redundant support, compiled depth, and nearly dependent tangent modes to accumulate \cite{Yao2021,Linteau2024,Zhang2025}.  Second, a suitable manifold can still be realized badly.  Ridge loading, eigencutoffs, damped inverses, and local subdivision are established numerical tools, but a condition number does not determine whether regularization should increase or decrease.  Stronger damping can suppress unstable modes or erase useful motion; a large physically accurate velocity instead calls for a smaller local step.  Structural expressivity miss must therefore be separated from numerical realization miss.

\subsection{AP-McLachlan: joint support-patch dynamics}
AP-McLachlan first treats structural maintenance as one bidirectional support-patch decision.  At a structural checkpoint it evaluates one generalized exchange patch $\mathcal B_k=(B_k^{\ominus},B_k^{\oplus})$.  Empty components recover stay, pure append, and pure prune; two nonempty components define a true exchange evaluated jointly on the patched support.  When the McLachlan distance exceeds its cut, the full generalized-exchange family is active and every patch is ordered first by its signed change in realized captured drift, with resource utility confined to a bounded numerical tiebreaker; after the cut is met, the addition component is held empty and the ordinary cost-aware prune utility resumes.  Prune retires support only after the realized patched solve, reduced-state refit, same-checkpoint ray and velocity continuity, cooldown, and atomic rejection authorize removal.  Among the many possible combinatorial support-search strategies, the implementation enumerates complete deletion-cardinality rungs under an explicit work budget, so all empty- and nonempty-component limits compete under the checkpoint's active ranking rule before any finalist is materialized.

After the support decision, the numerical lane determines how to realize Eq.~(\ref{eq:mclachlan_velocity}) on the accepted fixed support.  It compares stronger and weaker inverse policies, uses state-space motion and temporal kink to nominate local subdivision, selects the least-intervention admissible response, and returns to the base solve through release hysteresis.  Structural maintenance and numerical stabilization therefore act on different questions: which tangent directions should be active, and how should motion on that active tangent space be realized over the finite step?

The bounded contribution is the coordination of separate structural and realized numerical diagnostics with a checkpoint finalist family containing stay, append, continuity-certified online prune, and conditional true exchange.  We do not claim the invention of McLachlan propagation, regularized inversion, residual-triggered growth, or zero-angle append.  Within the gate-based adaptive variational-dynamics literature reviewed through July 18, 2026, we found no earlier method combining these elements as one measured checkpoint policy; this negative-priority statement is bounded by that review and will be updated as the remaining primary-source verification queue is completed.  Exact reference trajectories are excluded from every online decision and enter only the post-run comparisons.

\section{McLachlan checkpoint geometry and the two failure modes}
\label{sec:dynamics_method}

% Purpose-first method-opening contract: each substantive method subsection
% begins with the local failure, the construction used to address it, and the
% score, reduced model, controller object, or decision passed downstream.

A compact ground-state ansatz from ADAPT-style state preparation can be reused to simulate quench-driven real-time dynamics \cite{Grimsley2019,Tang2021,Yordanov2021}. McLachlan dynamics is NISQ-compatible under the constraint that the initialized ansatz circuit remains NISQ-compatible \cite{McLachlan1964,LiBenjamin2017VQS,Yuan2019VQS}. % \placeholder{Cites ADAPT initial-state preparation, the least-squares variational principle, and early near-term variational quantum simulation routes.}
As such, preparation of the initial state is not an auxiliary detail, but it is not the online decision law: the static ADAPT ansatz supplies the starting manifold on which the dynamic controller operates.

\subsection{Checkpoint geometry}
Exact Schr\"odinger motion generally leaves a finite variational manifold.  We therefore evaluate the projected tangent geometry and physical drift appearing in Eq.~(\ref{eq:mclachlan_velocity}) at each checkpoint.  These objects determine the parameter velocity used by the time integrator and by the downstream structural and numerical diagnostics.

Let time points be $0=t_0<\cdots<t_{N_t}=T_{\rm f}$, with step $\Delta t_k=t_{k+1}-t_k$. During propagation the Pauli labels are fixed and only the driven coefficients change:
\begin{equation}
H(t)=\sum_{\alpha\in\mathcal I_{\stat}}h_\alpha P_\alpha
+\sum_{\beta\in\mathcal I_{\drv}}c_\beta(t)P_\beta .
\label{eq:dynamics_sampling}
\end{equation}
We write $H(t_k)$ for the Hamiltonian evaluated at the checkpoint sample time.
For the inherited ADAPT ansatz, $\ket{\psi_k}=U(\theta_k)\ket{\psi_{\rm ref}}$, the parameterization supplies flattened coordinates $\theta_k\in\mathbb R^{N_k}$. Let $\Pi_k=I-\ket{\psi_k}\!\bra{\psi_k}$ be the horizontal projector.  The algorithm estimates, through measurement-compatible tangent observables, the horizontally projected tangent block
\begin{equation}
\bar T_k=\Pi_k T_k
=\Pi_k\left[
\partial_{\theta_{k,1}}\ket{\psi_k},\ldots,
\partial_{\theta_{k,N_k}}\ket{\psi_k}
\right],
\label{eq:projected_tangent_block}
\end{equation}
and the instantaneous projected Schrödinger drift
\begin{equation}
\bar b_k=-\ii\left(H(t_k)-\expval{H(t_k)}_{\psi_k}\right)\ket{\psi_k}.
\label{eq:projected_schrodinger_drift}
\end{equation}

The least-squares mismatch between the projected tangent velocity and the Schrödinger drift is the McLachlan residual on the prepared circuit manifold \cite{McLachlan1964,Yuan2019VQS,Barison2021}.
% \placeholder{Cites the variational least-squares residual and its parameterized-circuit dynamics form.}
\begin{equation}
\begin{aligned}
\mathfrak m_k(v)
&:=\|\bar T_kv-\bar b_k\|_2^2 \\
&=v^\top G_kv-2v^\top f_k+\|\bar b_k\|_2^2,
\end{aligned}
\label{eq:mclachlan_metric_objective}
\end{equation}
where
$G_k=\Re(\bar T_k^\dagger\bar T_k)$,
$f_k=\Re(\bar T_k^\dagger\bar b_k)$, and
$J_k:=\{1,\ldots,N_k\}$ is the active ansatz-coordinate support at time $t_k$.
Solving Eq.~(\ref{eq:mclachlan_velocity}) on the resolved tangent support gives $\dot\theta_k=G_k^\dagger f_k$.  The resulting velocity defines the classical IVP for $\theta(t)$ and hence the propagated state $\ket{\psi(\theta(t))}$.  Numerical rank selection, damping, and whitening concern the realization of this equation, not the McLachlan principle itself, and are collected in Appendix~\ref{app:technical_realization}.

\subsection{Structural and numerical failure}

A trajectory can fail because the active tangent support cannot represent the physical drift, or because the implemented inverse realizes an available tangent motion poorly.  We separate these mechanisms by comparing supported geometric explained drift with the realized stabilized motion.  This produces distinct structural and numerical misses for the checkpoint policy, rather than one residual that conflates support inadequacy with solve error.

For any coordinate subset $S\subseteq J_k$, define the supported explained drift
\begin{equation}
\Gamma_k(S)
=f_{k,S}^{\top}(G_{k,SS})^\dagger f_{k,S}.
\label{eq:explained_drift}
\end{equation}
$\Gamma_k$ is the idealized capacity of the support under an exact pseudo-inverse.  The structural scores below instead use the drift captured by the motion the solver actually returns.  For the realized velocity $v_S$ produced by the declared inverse policy on $S$ (retained-mode truncation, ridge, and damping included), the realized captured drift is
\begin{equation}
Q_k(S)
=2f_{k,S}^{\top}v_S-v_S^{\top}G_{k,SS}v_S.
\label{eq:realized_captured_drift}
\end{equation}
$Q_k(S)$ equals the reduction of the McLachlan objective in Eq.~(\ref{eq:mclachlan_metric_objective}) achieved by $v_S$, is exact under any regularized inverse, and reduces to $\Gamma_k(S)$ when $v_S=(G_{k,SS})^\dagger f_{k,S}$.  $\Gamma_k$ is retained as the idealized diagnostic; $Q_k$ is the scoring authority for every structural decision.
The normalized structural miss is
\begin{equation}
\rho_{\rm expr,k}(J_k)
=\operatorname{clip}_{[0,1]}\!\left(
\frac{\|\bar b_k\|_2^2-\Gamma_k(J_k)}
{\|\bar b_k\|_2^2+\eps_{\rm norm}}
\right).
\label{eq:active_structural_miss}
\end{equation}
Large $\rho_{\rm expr,k}$ means that resolved tangent support is absent.  By contrast, $\rho_{\rm num,k}$ below compares the realized stabilized velocity with the best supported least-squares motion on the same $J_k$.  A checkpoint can therefore realize an inadequate support accurately or realize an adequate support poorly.

At each checkpoint, the structural lane may choose one support patch, after which the numerical lane chooses how to realize motion on the accepted fixed support.  Both decisions use only current or previously acquired measurement-compatible data.  Exact or classically propagated reference trajectories are excluded.

\section{Joint support-patch maintenance}
\label{sec:structural_support}

Structural maintenance changes \(J_k\), the active ansatz-coordinate support, when the current manifold lacks useful directions or retains directions that are redundant, expensive, or harmful to the tangent solve.  Every structural candidate is evaluated against the same measured checkpoint geometry before the accepted support is passed to numerical realization.

The support scores use the hardware-cost factor \(K_{t_k}(B)\ge1\) introduced for static adaptive-ansatz construction in Ref.~\cite{Strobel2026StaticAdapt}.  This \(K_{t_k}\) is a compiled-resource weight, not the tangent Gram matrix \(G_k\).  Geometric gain or loss remains the physical authority; cost only ranks otherwise admissible patches.

\subsection{Generalized exchange support patch}

When $L_k^2$ exceeds its declared cut, the structural action is always one generalized exchange patch; append, prune, and true exchange are not selected by separate upstream rules.  The patch is
\begin{equation}
\mathcal B_k=(B_k^{\ominus},B_k^{\oplus}),
\qquad
J_k^{\mathcal B}
=
(J_k\setminus I_{B,k}^{\ominus})
\cup I_{B,k}^{\oplus}.
\label{eq:support_patch}
\end{equation}
The action is classified only by which components of the selected patch are empty: $B_k^{\ominus}=B_k^{\oplus}=\varnothing$ gives stay, $B_k^{\ominus}=\varnothing$ gives the pure-append limit, $B_k^{\oplus}=\varnothing$ gives the pure-prune limit, and two nonempty components give true exchange.  Thus \emph{generalized exchange} names the complete two-component action, whereas \emph{true exchange} is reserved for a patch that both removes and adds support.  A true exchange recomputes the append gain against the post-delete support, the deletion loss against the patched support, and the McLachlan solve on \(J_k^{\mathcal B}\).  It then applies the same reduced-state and velocity-continuity tests as the pure-prune limit.  Failure of a true exchange does not authorize either one-sided limit; any patch with an empty component must appear and pass independently as a finalist.

\subsection{Pure append limit}

Set \(B_k^{\ominus}=\varnothing\).  Let \(B_k^{\oplus}\) be a candidate append batch and \(I_{B,k}^{\oplus}\) its new coordinate indices.  Added coordinates enter at zero amplitude, so the checkpoint state is unchanged, and this invariance holds at every circuit position.  Each candidate therefore carries an insertion position: the new rotation may be placed at any cut of the existing product circuit, positions related by block-swap commutation certificates generate identical one-parameter families, and only one representative per equivalence class is enumerated.  Distinct retained cuts produce distinct tangent directions and are scored as distinct candidates.  The normalized gain from refitting the old and new tangent coordinates together is
\begin{equation}
\Delta_k^{\oplus}(B_k^{\oplus})
=
\frac{
\left[
Q_k(J_k\cup I_{B,k}^{\oplus})
-Q_k(J_k)
\right]_+
}{
\|\bar b_k\|_2^2+\eps_{\rm norm}
},
\label{eq:append_gain}
\end{equation}
and the resource-weighted ranking score is
\begin{equation}
S_k^{\oplus}(B_k^{\oplus})
=
\frac{
\Delta_k^{\oplus}(B_k^{\oplus})
}{
K_{t_k}(B_k^{\oplus})^{\alpha_{\rm append}}
}.
\label{eq:append_score}
\end{equation}
The batch is evaluated jointly rather than by summing singleton gains.

Gain alone is insufficient when the incoming tangent block duplicates the active span; duplicate Pauli words generate the same one-parameter family and are removed by operator identity before scoring, and an append batch must yield a finite augmented McLachlan solve within the declared conditioning bound.  Multi-child insertion batches are ordered words over the inserted children and the survivors, quotiented to a canonical normal form by pairwise commutation so that only inequivalent orderings are enumerated; wider insertion frontiers are acquired only while smaller families certify no patch and the structural-repair predicate remains active.  Whitening and retained-mode compression are implementation accelerations collected in Appendix~\ref{app:technical_realization}; they do not replace Eq.~(\ref{eq:append_gain}).

Append scoring requires overlaps involving candidate tangent directions and can therefore require additional quantum measurements beyond the active-support McLachlan data.  After those candidate blocks have been measured, batch assembly, Schur projection, ranking, and the augmented solve are classical linear algebra.

\subsection{Pure prune limit}

Set \(B_k^{\oplus}=\varnothing\).  For a candidate deletion batch \(B_k^{\ominus}\) with active coordinate set \(I_{B,k}^{\ominus}\subseteq J_k\), the normalized loss after all surviving tangent coordinates refit is
\begin{equation}
\mathcal L_k^{\ominus}(B_k^{\ominus})
=
\frac{
\left[
Q_k(J_k)
-Q_k(J_k\setminus I_{B,k}^{\ominus})
\right]_+
}{
\|\bar b_k\|_2^2+\eps_{\rm norm}
}.
\label{eq:prune_loss}
\end{equation}
Equation~(\ref{eq:prune_loss}) is the backward counterpart of Eq.~(\ref{eq:append_gain}): append asks how much explained drift is gained by enlarging the support, whereas prune asks how much is lost after deleting a block and refitting the survivors.  The positive-part operation makes $\mathcal L_k^{\ominus}$ a damage quantity: it is positive for a harmful deletion and zero for both a harmless deletion and a deletion that improves the realized solve.  Low loss can therefore nominate a block for cost-aware compression, but it cannot order deletion-containing patches while the checkpoint owes accuracy.  The signed net drift change defined below has that authority; resource cost, historical redundancy, Schur degeneracy, and conditioning relief remain secondary nomination information.

The quantity in Eq.~(\ref{eq:prune_loss}) has an exact algebraic form, and it is worth stating because it fixes what pruning can and cannot do.  Partition the checkpoint Gram system between the coordinates a candidate retains, \(R=J_k\setminus D\), and those it proposes to delete, \(D\),
\begin{equation}
G_k=
\begin{bmatrix} G_{RR} & G_{RD}\\ G_{DR} & G_{DD}\end{bmatrix},
\qquad
f_k=\begin{bmatrix} f_R\\ f_D\end{bmatrix},
\end{equation}
and define the reverse Schur novelty block and its drive residual
\begin{equation}
\begin{aligned}
S_{D|R}&=G_{DD}-G_{DR}G_{RR}^{+}G_{RD},\\
r_{D|R}&=f_D-G_{DR}G_{RR}^{+}f_R .
\end{aligned}
\label{eq:schur_novelty}
\end{equation}
Then under an exact, unregularized solve
\begin{equation}
Q_k(J_k)-Q_k(R)=r_{D|R}^{\mathsf T}S_{D|R}^{+}r_{D|R}\;\ge\;0 .
\label{eq:schur_identity}
\end{equation}
The identity requires \(f_k\in\operatorname{range}(G_k)\), which this route satisfies by construction rather than by assumption: \(G_k=T_k^{\mathsf T}T_k\) and \(f_k=T_k^{\mathsf T}\bar b_k\) are formed from the same tangent matrix, so \(f_k\in\operatorname{range}(T_k^{\mathsf T})=\operatorname{range}(G_k)\), and restricting to a coordinate subset restricts \(T_k\) to those columns and preserves the property.

Two consequences follow, and the route is built around them.  First, the right-hand side of Eq.~(\ref{eq:schur_identity}) is a quadratic form in a positive-semidefinite Schur complement, so \emph{under exact projection a deletion can never increase captured drift}, and therefore can never lower \(L_k^2=2[\|\bar b_k\|_2^2-Q_k]_+\).  Any helpful deletion observed in a run is possible only because the realized solve is not exact --- ridge regularization, rank truncation, damping, or the trust-region refit.  Pruning helps by relieving the conditioning of the realized solve, not by improving the variational manifold, and this is the precise content of the accuracy half of the prune claim.  Second, \(S_{D|R}\) measures how much of the deleted tangent block cannot be synthesized from the retained tangents and \(r_{D|R}\) how much forcing remains aligned with that novel part, so Schur degeneracy \emph{nominates} a coordinate but cannot \emph{authorize} its removal: only the realized post-deletion solve establishes that removing it helps.  Eq.~(\ref{eq:schur_identity}) is pinned numerically in the implementation, on random systems and on the route's own construction, so that a change to the geometry, the inverse policy, or the scoring cannot silently break it.

Because Eq.~(\ref{eq:prune_loss}) uses submatrices of the active \(G_k\) and \(f_k\) already acquired for McLachlan propagation, prune nomination itself requires no new quantum measurements.  A deletion commit is stricter.  The candidate block is first set to the identity and, if necessary, the surviving parameters are locally refit to the pre-patch state by a bounded trust-region minimization of the Fubini--Study infidelity to the checkpoint ray: the coordinate word is held fixed, the angles move inside a declared box radius, the refit is accepted only if it strictly reduces the infidelity, and finalists already on the ray (every pure zero-angle insertion) are left untouched.  This reduced-state fit or overlap certification can require additional circuits on a QPU.

Let \(J_k^{\mathcal B}\) be the proposed post-patch support and \(\vartheta_{\mathcal B,k}^{\star}\) the zero-transported or locally refit parameters on that support.  Deletion can commit only if
\begin{equation}
d_{\rm FS}\!\left(
\psi_k,
\psi_k^{\mathcal B}
\right)
\le\tau_{\rm patch},
\qquad
\eta_{{\rm patch},k}^{\psi}(\mathcal B_k)
\le\tau_{\rm patch}^{\rm vel},
\label{eq:patch_continuity}
\end{equation}
where
\begin{equation}
\eta_{{\rm patch},k}^{\psi}(\mathcal B_k)
=
\frac{
\left\|
\bar T_k^{\mathcal B}\dot\vartheta_{\mathcal B,k}^{\star}
-\bar T_k\dot\theta_k
\right\|_2^2
}{
\|\bar b_k\|_2^2+\eps_{\rm norm}
}.
\label{eq:patch_velocity_continuity}
\end{equation}
Persistence prevents a transiently redundant block from being removed after one checkpoint; cooldown and atomic finalist rejection prevent repeated unsafe attempts.  The history and conditioning nomination score is given in Appendix~\ref{app:technical_realization}.

\subsection{Operating decision rule}

Collecting the pieces, the rule evaluated at every checkpoint is compact.  Deletions are always eligible: a candidate deletion is a row and column selection of geometry already acquired for propagation, so its structural evaluation costs no additional measurement.  Below the McLachlan-distance cut the addition component is constrained to $B_k^{\oplus}=\varnothing$, leaving stay and the pure-prune limit.  Above the cut, candidate tangent measurements open the full generalized-exchange family, with either component allowed to be empty.  The controller therefore ranks complete patches and assigns the labels append, prune, or true exchange only from the empty-component pattern of the selected action.  For every admitted patch define the signed normalized change in realized captured drift
\begin{equation}
\delta_k(\mathcal B)
=
\frac{Q_k(J_k^{\mathcal B})-Q_k(J_k)}
{\|\bar b_k\|_2^2+\eps_{\rm norm}}.
\label{eq:signed_patch_drift}
\end{equation}
At the frozen checkpoint geometry, positive $\delta_k$ is equivalent, up to the nonnegative numerical clip, to a reduction of
$L_k^2=2[\|\bar b_k\|_2^2-Q_k]_+$; explicitly,
$L_k^2(J_k)-L_k^2(J_k^{\mathcal B})=2(\|\bar b_k\|_2^2+\eps_{\rm norm})\delta_k(\mathcal B)$ whenever the clip is inactive.  The ordinary resource-aware utility is
\begin{equation}
\begin{aligned}
U_k(\mathcal B)
&=U_k^{\oplus}(\mathcal B)+U_k^{\ominus}(\mathcal B)
+w_\delta\delta_k(\mathcal B),\\
U_k^{\oplus}(\mathcal B)
&=\frac{\Delta_k^{\oplus}}{K_{t_k}^{\alpha_{\rm ins}}},
\qquad
U_k^{\ominus}(\mathcal B)
=\frac{K_{t_k}^{\alpha_{\rm del}}}{\mathcal L_k^{\ominus}+\eps_L}.
\end{aligned}
\label{eq:operating_score}
\end{equation}
where an absent insertion or deletion component contributes zero.  While $L_k^2$ exceeds the cut, the operating key is instead
\begin{equation}
\begin{aligned}
\mathbf R_k^{\rm debt}(\mathcal B)
&=\left(R_{\delta,k}(\mathcal B),R_{U,k}(\mathcal B)\right),\\
R_{\delta,k}(\mathcal B)
&=\operatorname{round}\!\left[\frac{\delta_k(\mathcal B)}{\tau_\delta}\right],
\qquad \tau_\delta=10^{-12},\\
R_{U,k}(\mathcal B)
&=\tanh\!\left[U_k^{\oplus}(\mathcal B)+U_k^{\ominus}(\mathcal B)\right].
\end{aligned}
\label{eq:debt_score}
\end{equation}
where $U_k^{\oplus}$ and $U_k^{\ominus}$ are the first two terms of Eq.~(\ref{eq:operating_score}).  Candidates are ordered lexicographically in descending $\mathbf R_k^{\rm debt}$: the first component places signed realized-drift changes into declared accuracy bins of width $\tau_\delta$, and the bounded resource utility orders candidates only within one bin.  The reciprocal deletion utility therefore cannot overturn a resolved accuracy ordering.  Once $L_k^2$ is below the cut, candidates are ranked by Eq.~(\ref{eq:operating_score}) and the route performs cost-aware support maintenance.  Ranked candidates are certified one at a time against the ray and velocity-continuity tests of Eqs.~(\ref{eq:patch_continuity})--(\ref{eq:patch_velocity_continuity}) after a bounded local refit, and the first patch that passes commits atomically.  During debt, the materialized candidate must also strictly lower the realized $L_k^2$ before it can commit.  A finiteness and conditioning assertion protects every finalist but is not a tuned threshold: across the reported runs the largest committed retained condition number is five orders of magnitude below its bound.  The history and conditioning extensions of Appendix~\ref{app:technical_realization} are carried at zero weight in the provisional runs reported below.

\subsection{Bounded support-patch search used here}

The ideal support-patch problem is combinatorial.  Exhaustive enumeration, mixed-integer optimization, branch-and-bound, greedy forward--backward selection, and local exchange are all possible search strategies for related subset problems \cite{Natarajan1995SparseApprox,BertsimasKingMazumder2016,DasKempe2011,Elenberg2018WeakSubmodularity}.  AP-McLachlan enumerates complete candidate families under one explicit work budget because candidate tangent measurements, not only classical subset enumeration, contribute to the checkpoint cost.

Deletion candidates are organized into cardinality rungs: the rung at cardinality $d$ contains every hard-feasible deletion of $d$ active coordinates.  A rung is admitted whole or not at all against a joint-work budget on scored candidates, and every admitted rung is enumerated completely; nothing inside a rung is sampled or prefiltered by score.  For each admitted deletion branch $D$, the deletion-conditioned insertion candidates are the deduplicated pool children placed at their commutation-reduced cuts remapped onto the post-deletion survivor word.  The resulting family contains generalized exchange patches with both empty and nonempty components, each evaluated on its complete patched support; append, prune, and true exchange are outcome labels determined after selection.  Equation~(\ref{eq:debt_score}) orders the full family by signed realized drift while accuracy debt remains; Eq.~(\ref{eq:operating_score}) orders the $B_k^{\oplus}=\varnothing$ face by deletion loss, history, conditioning, and resource cost after the distance cut is met.

Certification then proceeds one finalist at a time in descending structural score.  A finalist is materialized (delete, remap cuts, insert at zero angle), locally refit toward the checkpoint ray, and passed through the hard commit gates: a finite patched solve within the conditioning bound for every finalist, and the ray-continuity and velocity-smoothness tests of Eqs.~(\ref{eq:patch_continuity})--(\ref{eq:patch_velocity_continuity}) for deletion-containing patches.  The first finalist that passes commits atomically; a failed finalist is rejected whole and certification advances.  Two admission bounds keep certification work finite without touching the ranking: a per-level budget on materialized finalists, and a per-branch budget that skips further insertion variants of a deletion branch after repeated failures, since branch-level score terms are shared by every variant of that branch.  Wider multi-child insertion frontiers are acquired only while the structural-repair predicate remains active and the work budget admits the complete family.

Measurement economics fixes which face of the generalized-exchange family is searched.  Insertion candidates require new tangent-overlap measurements, so $L_k^2$ above its declared cut opens the complete two-component family.  Pure-prune-limit patches are row and column selections of geometry already acquired for propagation and are considered at every checkpoint; below the cut the search is therefore restricted to the measurement-free $B_k^{\oplus}=\varnothing$ face.  Exact trajectories do not enter the enumeration, ranking, or commit tests.  Further history, conditioning, and retry conventions are collected in Appendix~\ref{app:technical_realization}.

\section{Fixed-support numerical stabilization}
\label{sec:numerical_solve_repair}
\label{sec:numerical_stabilization}

An expressive tangent support can still fail numerically.  The inverse may lose motion that lies in the resolved tangent span, a finite time step may move too far for the local linearization, or the accepted state-space velocity may change discontinuously between neighboring time points.  These are three different failures, and the three diagnostics below measure them separately.

The three diagnostics are measured at every checkpoint, but they do not all open the same response.  Local subdivision is what repeatedly changes a step: across the trajectories reported here it opened $372$ times on excessive projected state displacement and $254$ times on the temporal-kink diagnostic.  The inverse-candidate search is what does not act: candidate conventions other than the base regularization were selected on two integration steps out of $885$, and six of $1.15\times10^{4}$ evaluated candidates were accepted (Appendix~\ref{app:diagnostic_only}).  The operating rule is therefore: subdivide the interval when the state-space step of Eq.~(\ref{eq:state_space_step}) or the temporal kink of Eq.~(\ref{eq:temporal_kink}) exceeds its bound; the numerical-miss channel opens no search in the reported configuration.

For any realized candidate velocity \(v\), the best geometric residual available on the current support is
\begin{equation}
\epsilon_{{\rm geo},k}^{2}
=
\left[
\|\bar b_k\|_2^2-\Gamma_k(J_k)
\right]_+,
\label{eq:geometric_residual}
\end{equation}
whereas its realized residual is
\begin{equation}
\epsilon_{{\rm real},k}^{2}(v)
=v^\top G_kv-2f_k^\top v+\|\bar b_k\|_2^2.
\label{eq:realized_residual}
\end{equation}
The normalized numerical miss is therefore
\begin{equation}
\rho_{{\rm num},k}(v)
=
\frac{
\left[\epsilon_{{\rm real},k}^{2}(v)
-\epsilon_{{\rm geo},k}^{2}\right]_+
}{
\|\bar b_k\|_2^2+\eps_{\rm norm}
}.
\label{eq:numerical_miss}
\end{equation}
Thus \(\rho_{{\rm num},k}\) measures motion lost by the realized solve relative to what the same tangent support can represent.  It is not an expressivity residual and does not decide whether a generator should be appended or removed.

The second diagnostic is the size of the proposed state-space step.  We use the larger of the tangent-plane estimate and the realized trial-state ray displacement,
\begin{equation}
\begin{aligned}
\delta_{{\rm step},k}(v)
&=\max\Bigl\{
\Delta t_k\sqrt{v^\top G_kv},
\\[-2pt]
&\hspace{3.3em}
d_{\rm FS}\!\left(
\psi(\theta_k),
\psi(\theta_k+\Delta t_kv)
\right)
\Bigr\}.
\end{aligned}
\label{eq:state_space_step}
\end{equation}
The third diagnostic detects a temporal kink in the physical tangent velocity,
\begin{equation}
\eta_{{\rm time},k}^{\psi}(v)
=
\frac{
\left\|
\bar T_kv-
\bar T_{k-1}\dot\theta_{k-1}^{\rm acc}
\right\|_2^2
}{
\tfrac12\left(
\|\bar b_k\|_2^2+
\|\bar b_{k-1}\|_2^2
\right)+\eps_{\rm norm}
}.
\label{eq:temporal_kink}
\end{equation}
For \(k=0\), the temporal-kink diagnostic is set to zero.




\paragraph*{Operating rule.}  The operating response is local subdivision: when \(\delta_{{\rm step},k}\) of Eq.~(\ref{eq:state_space_step}) or \(\eta_{{\rm time},k}^{\psi}\) of Eq.~(\ref{eq:temporal_kink}) exceeds its bound, the interval is subdivided, because damping is not the first cure for a physically accurate velocity evaluated over too large a step, nor for a velocity that has turned sharply between neighboring checkpoints.  Subdivision refines the internal integration grid only, leaving the reporting grid and the accepted support unchanged.  Hard invalidity --- nonfinite or dimensionally inconsistent checkpoint data --- rejects the step outright.  The bidirectional inverse search over candidate regularizations, together with the release rule that returns a held repair to the base convention, is specified in Appendix~\ref{app:diagnostic_only} with the measurements motivating its exclusion from the operating configuration.

The measurement boundary is important.  Once \(G_k\), \(f_k\), and \(\|\bar b_k\|_2^2\) have been estimated at a checkpoint, Eqs.~(\ref{eq:realized_residual}) and the tangent-plane part of Eq.~(\ref{eq:state_space_step}) are classical quadratic forms.  Testing stronger or weaker inverse policies on those same data requires no new quantum measurements.  The trial-state ray check requires an additional overlap experiment on a QPU when that stricter guard is enabled; in the present statevector diagnostics it is evaluated classically.  Local subdivision is different: it introduces new intermediate states and therefore requires new checkpoint evaluations.  Temporal continuity uses current and previously retained checkpoint data and never an exact reference trajectory.



\subsection{Time integration convention}

The accepted velocity defines the parameter IVP integrated over each reporting interval.  Euler and fourth-order Runge--Kutta are available diagnostic integrators; local subdivision changes only the internal integration grid, while observables remain reported on the declared common time grid.  Structural support changes occur only at declared checkpoints.

\section{Results}
\label{sec:results}
\label{sec:benchmarks}

The results identify mechanisms before surveying Hamiltonian regimes.  Every paired comparison holds fixed the initial state, static ansatz seed, Hamiltonian, drive, time grid, nominal inverse, observables, and compilation convention.  Exact trajectories are attached only after propagation to define reporting errors.

\subsection{Mechanism-level ablations}

The structural comparison holds numerical realization fixed and contrasts policy arms of one selector: append-only (deletions never eligible), prune-only (empty insertion pool), and the full exchange policy, each with and without the finalist trust-region refit, under one seed, Hamiltonian, drive, grid, and inverse policy.  It reports structural miss, accepted patch types by kind, conditional gain and loss, same-checkpoint continuity, final support, compiled cost, and the cumulative count of insertion candidates whose tangent overlaps had to be measured --- the measurement ledger that separates gated insertion enumeration from measurement-free pruning.  The numerical-repair comparison then holds the support fixed and contrasts passive telemetry under the nominal solve with the intervening state-space policy.  It reports numerical miss, state-space displacement, temporal kink, inverse-policy changes, local subdivisions, unsupported checkpoints, and post-run energy and observable errors.  These comparisons separate support maintenance from numerical realization before their interaction is interpreted.

\begin{table}[t]
\caption{Mechanism-level evidence plan.  Numerical entries remain provisional until transferred from source-locked paired trajectories under one seed, drive, grid, observable set, and compilation contract.  Exact-reference errors are post-run quantities.}
\label{tab:mechanism_ablation}
\begin{ruledtabular}
\begin{tabular}{lcccc}
Configuration & $\bar\epsilon_E$ & $\epsilon_D$ & $\int\rho_{\rm expr}$ & $D_{\twq}^{\rm final}$\\
\colrule
fixed support, nominal solve & \tentative{--} & \tentative{--} & \tentative{--} & \tentative{--}\\
fixed support, solve repair & \tentative{--} & \tentative{--} & \tentative{--} & \tentative{--}\\
append-only & \tentative{--} & \tentative{--} & \tentative{--} & \tentative{--}\\
exchange, no refit & \tentative{--} & \tentative{--} & \tentative{--} & \tentative{--}\\
exchange with refit & \tentative{--} & \tentative{--} & \tentative{--} & \tentative{--}\\
\end{tabular}
\end{ruledtabular}
\end{table}

\subsection{Support-removal evidence at production horizon}
\label{sec:production_matrix}

The claim this paper makes is that online support removal reduces circuit cost
at matched or better trajectory accuracy for ans\"atze that carry redundant
support.  Table~\ref{tab:production_matrix} tests it on three starting points
spanning the redundancy axis, each propagated to $t=5$ over $126$ reporting
intervals under the same drive, integrator, inverse policy, repair
configuration, candidate pool, and computational guards.  Each cell records a
physics fingerprint binding its seed, cutoff, drive, and grid; the arms within
a seed group were verified to share that fingerprint before being tabulated, so
comparability is a property of the artifacts rather than of the launch
procedure.

\begin{table}[t]
\caption{Structural policies at production horizon ($t=5$, $126$ intervals, driven).  Errors are against the exact time-ordered trajectory of the same truncated Hamiltonian.  ``Measured'' counts candidate directions whose tangent overlaps had to be estimated, the quantum cost of structural maintenance; deletions contribute nothing to it because they reuse geometry already acquired for propagation.  \tentative{Preliminary: this horizon is short relative to the intended regime, and continuation runs from these $t=5$ end states are planned.}}
\label{tab:production_matrix}
\begin{ruledtabular}
\begin{tabular}{llcccc}
Seed & Policy & $\overline{|\Delta E|}$ & $|\Delta E|_{t=T}$ & $N_{\rm param}$ & measured\\
\colrule
\multirow{3}{*}{redundant} & append-only & \tentative{$1.6\times10^{-2}$} & \tentative{$3.1\times10^{-2}$} & \tentative{668} & \tentative{384}\\
 & adaptive append & \tentative{$1.8\times10^{-2}$} & \tentative{$4.5\times10^{-2}$} & \tentative{866} & \tentative{1000}\\
 & exchange & \tentative{$1.9\times10^{-2}$} & \tentative{$2.4\times10^{-2}$} & \tentative{606} & \tentative{400}\\
\colrule
\multirow{3}{*}{fermionic} & append-only & \tentative{$1.49\times10^{-2}$} & \tentative{$2.42\times10^{-2}$} & \tentative{20} & \tentative{0}\\
 & adaptive append & \tentative{$1.49\times10^{-2}$} & \tentative{$2.42\times10^{-2}$} & \tentative{20} & \tentative{0}\\
 & exchange & \tentative{$1.49\times10^{-2}$} & \tentative{$2.42\times10^{-2}$} & \tentative{10} & \tentative{0}\\
\colrule
\multirow{3}{*}{compact} & append-only & \tentative{$5.7\times10^{-3}$} & \tentative{$9.1\times10^{-3}$} & \tentative{77} & \tentative{392}\\
 & adaptive append & \tentative{$3.9\times10^{-2}$} & \tentative{$6.0\times10^{-2}$} & \tentative{233} & \tentative{1000}\\
 & exchange & \tentative{$1.7\times10^{-2}$} & \tentative{$2.4\times10^{-2}$} & \tentative{13} & \tentative{792}\\
\end{tabular}
\end{ruledtabular}
\end{table}

On the redundant seed --- a fixed-structure VQE ansatz of $620$ coordinates,
whose redundancy is a property of that ansatz family rather than an arrangement
made for this test --- the three policies reach the same mean error to within
their spread, while exchange ends on the smallest support and with the smallest
final error.  Growth without removal costs both: the adaptive-append comparator
ends $43\%$ larger than exchange and with roughly twice its final error, having
spent two and a half times as many measured candidate directions.

The fermionic control isolates the removal mechanism from the phonon sector.
There no policy adds a coordinate, so append-only and adaptive append are
indistinguishable; exchange reaches the same trajectory to four significant
figures on half the support, using no additional measurements at all.  This is
the claim in its simplest form: the removed coordinates were carrying no
trajectory information, and removing them is free.

The compact seed marks the boundary.  With little redundancy to remove,
exchange compresses hardest ($28\to13$) but pays in accuracy against
append-only, and the trade rather than the win is what the data supports.  We
report it because the claim is scoped to redundancy-carrying support, and the
boundary of a claim belongs in evidence rather than in a caveat.


\subsection{Observable tracking on the driven Holstein seed}
\label{sec:hh_observables}

Energy error summarizes accuracy but hides whether a policy tracks the physics
or merely its scale.  The observable evidence is reported with the
adaptive-append comparison in Sec.~\ref{sec:comparator}, on the drive that
stresses the structural rule hardest; the doublon fraction and the site
occupations are followed to $t=10$ against the exact time-ordered trajectory of
the same truncated Hamiltonian.

\subsection{How comparisons are constructed}
\label{sec:comparison_protocol}

Every comparison reported below is taken \emph{at fixed accuracy}, not at a
fixed activation threshold, and this choice is forced by measurement rather than
preference.  Each adaptive method carries its own McLachlan-distance cut
\(L^2_{\rm cut}\) below which it stops enlarging the ansatz, and the natural
comparison is to set both to a common value.  On this problem that is not sound.
Sweeping the cut over two decades on three drives, tightening it improves
accuracy monotonically on one of three drives for the comparator and none of
three for the route of this work, while a hundredfold change of the cut moves
the support by roughly ten coordinates; on the weak slow drive the two methods
return identical trajectories at the two loosest cuts because neither ever
triggers.  The cut therefore has a dead zone, a live zone, and non-monotone
behavior inside the live zone, and a method's quality cannot be read off one
setting of its own dial.

We instead fix an accuracy target \(\varepsilon\) and ask what each method costs
to reach it.  For every cell we run a common ladder of cuts, record the mean
absolute energy error against the exact time-ordered trajectory and the final
support size at each rung, and report the smallest support among rungs meeting
\(\varepsilon\); that support size is the comparable cost.  A cell that no rung
reaches is reported as not converged rather than averaged over, and a target is
only meaningful when both methods receive the same ladder and the same selection
rule.  Because selecting the best rung is a minimum over a noisy ladder, we
report the tightest-rung comparison --- which involves no selection --- beside
the best-rung one, and do not quote best-rung ratios as headline figures.

Three axes are held identical across any pair of arms being compared: the
physics (seed, drive, reporting grid, candidate pool), the inner numerical
method (integrator, Tikhonov regularization, damping, pseudo-inverse cutoff),
and the step-size control law.  The last is a genuine axis rather than a
property of a method: both methods bound their step and differ in the quantity
bounded, this route bounding state displacement and subdividing while the
comparator bounds \(\max_\mu|\delta\theta_\mu|\), and the integrator matters
independently --- fourth-order integration outperforms first-order by roughly a
factor of three and a half for \emph{both} methods, so arms differing in
integrator are not comparable.  The full axis specification, the locked values,
and the per-cell reporting contract are given in the lane protocol document
\texttt{agent\_guidance/time-dynamics/paper-ii-comparison-protocol.md}, and the
exact command and complete parameter set for every reported run are generated
into \texttt{agent\_guidance/time-dynamics/paper-ii-run-parameters.md}.

Under this protocol the growth-only ablation of this route reaches
\(\varepsilon=10^{-4}\) on three of three drives against the comparator's two of
three, at equal or smaller support (\tentative{47} against \tentative{50}, and
\tentative{44} against \tentative{47}).  \tentative{The corresponding ladder for
the full exchange route is in progress; whether pruning attains a lower
converged error, or the same error at smaller support, is the measurement that
decides the central claim of this work, and no fixed-accuracy statement about
pruning is made here in its absence.}

\subsection{Frozen adaptive-append baseline}
\label{sec:competitor_baseline}

The comparator is measured once, across the full drive set, and then held fixed;
every subsequent result for the route of this work is reported against that
frozen baseline rather than against a comparator re-run alongside it.  Two
variants are recorded because they answer different questions.  The
\emph{published} variant reproduces Ref.~\cite{Yao2021} as specified --- first-order
integration, Tikhonov regularization $\xi=10^{-6}$, and the step
dynamically adjusted so that $\max_\mu|\delta\theta_\mu|\le\delta\theta_{\max}=5\times10^{-3}$,
which is that method's own stated stabilization mechanism.  The
\emph{shared-numerics} variant runs the same decision rule on this route's inner
numerical method, isolating the structural rule from the numerics.  Neither is a
substitute for the other: a claim about the published method must cite the first,
and a claim about structural policy must cite the second.

\begin{table*}[t]
\caption{Frozen adaptive-append baseline on the driven Hubbard--Holstein seed ($L=2$, $n_{\rm ph}^{\max}=1$) to $t=10$ over $251$ checkpoints, against the exact time-ordered trajectory of the same truncated Hamiltonian.  Both variants share the seed, drive, reporting grid, and the full deduplicated candidate pool of $125$ words, and neither is capped: each appends until its McLachlan-distance cut is met.  \tentative{Provisional pending the source-locked matched matrix for this work's route.}}
\label{tab:competitor_baseline}
\begin{ruledtabular}
\begin{tabular}{lcccccc}
& \multicolumn{3}{c}{published numerics} & \multicolumn{3}{c}{shared numerics}\\
\cline{2-4}\cline{5-7}
Drive & $\overline{|\Delta E|}$ & $\overline{|\Delta n_{\rm d}|}$ & $N_{\rm param}$ & $\overline{|\Delta E|}$ & $\overline{|\Delta n_{\rm d}|}$ & $N_{\rm param}$\\
\colrule
strong fast ($A{=}2.4$, $\omega{=}3$) & \tentative{$6.97\times10^{-3}$} & \tentative{$3.24\times10^{-3}$} & \tentative{47} & \tentative{$3.84\times10^{-3}$} & \tentative{$1.01\times10^{-3}$} & \tentative{47}\\
fast weak ($A{=}0.6$, $\omega{=}3$) & \tentative{$7.79\times10^{-3}$} & \tentative{$2.10\times10^{-3}$} & \tentative{48} & \tentative{$2.28\times10^{-4}$} & \tentative{$7.45\times10^{-4}$} & \tentative{48}\\
slow strong ($A{=}1.2$, $\omega{=}1$) & \tentative{$8.74\times10^{-3}$} & \tentative{$3.11\times10^{-3}$} & \tentative{48} & \tentative{$2.44\times10^{-4}$} & \tentative{$1.01\times10^{-3}$} & \tentative{51}\\
mid resonance ($A{=}1.2$, $\omega{=}2$) & \tentative{$3.10\times10^{-2}$} & \tentative{$5.08\times10^{-3}$} & \tentative{64} & \tentative{$2.06\times10^{-4}$} & \tentative{$8.94\times10^{-4}$} & \tentative{52}\\
weak slow ($A{=}0.3$, $\omega{=}1$) & \tentative{$1.14\times10^{-3}$} & \tentative{$5.54\times10^{-4}$} & \tentative{47} & \tentative{$1.82\times10^{-4}$} & \tentative{$4.94\times10^{-4}$} & \tentative{48}\\
fast fast ($A{=}0.6$, $\omega{=}6$) & \tentative{$2.06\times10^{-2}$} & \tentative{$3.69\times10^{-3}$} & \tentative{47} & \tentative{$3.10\times10^{-3}$} & \tentative{$1.04\times10^{-3}$} & \tentative{50}\\
\end{tabular}
\end{ruledtabular}
\end{table*}

The two columns differ by up to two orders of magnitude on the same decision
rule --- most sharply under the fast weak drive, \tentative{$7.8\times10^{-3}$}
against \tentative{$2.3\times10^{-4}$} --- with support sizes that agree to
within a few coordinates throughout.  The comparator's accuracy on this problem
is therefore set largely by its numerical method rather than by its structural
rule, which is the reason both variants are carried: reporting only the
shared-numerics column would overstate the published method, and reporting only
the published column would attribute to structure what belongs to integration.


\par\addvspace{1.0ex}

\par\addvspace{1.0ex}

\subsection{Matched comparison at the frozen baseline}
\label{sec:matched_comparison}

Table~\ref{tab:matched_comparison} places this route against the frozen
adaptive-append baseline of Sec.~\ref{sec:competitor_baseline} under one inner
numerical method: fourth-order integration, Tikhonov regularization
$\xi=10^{-7}$, the same seed, drive set, reporting grid, candidate pool, and
computational guards.  The comparator is inexpensive relative to this route, so
it is re-measured whenever the shared numerics change rather than carried
forward from an earlier setting; a matched claim is only meaningful if the match
is re-established each time.

\begin{table}[t]
\caption{This route against the frozen adaptive-append baseline at matched numerics, driven Hubbard--Holstein seed ($L=2$, $n_{\rm ph}^{\max}=1$) to $t=10$.  Errors are against the exact time-ordered trajectory of the same truncated Hamiltonian.  \tentative{One seed and six drives; the deletion-under-debt policy is held at its default, and the remaining cell is in progress.}}
\label{tab:matched_comparison}
\begin{ruledtabular}
\begin{tabular}{lcccccccc}
& \multicolumn{3}{c}{adaptive append} & \multicolumn{3}{c}{this work} &\\
\cline{2-4}\cline{5-7}
Drive & $\overline{|\Delta E|}$ & $\overline{|\Delta n_{\rm d}|}$ & $N_{\rm param}$ & $\overline{|\Delta E|}$ & $\overline{|\Delta n_{\rm d}|}$ & $N_{\rm param}$ & $\Delta N$\\
\colrule
strong fast & \tentative{$9.12\times10^{-3}$} & \tentative{$2.73\times10^{-3}$} & \tentative{46} & \tentative{$3.10\times10^{-2}$} & \tentative{$8.56\times10^{-3}$} & \tentative{36} & \tentative{22\%}\\
fast weak & \tentative{$5.58\times10^{-4}$} & \tentative{$6.95\times10^{-4}$} & \tentative{46} & \tentative{$1.09\times10^{-2}$} & \tentative{$2.59\times10^{-3}$} & \tentative{34} & \tentative{26\%}\\
slow strong & \tentative{$8.00\times10^{-4}$} & \tentative{$1.37\times10^{-3}$} & \tentative{47} & \tentative{$3.21\times10^{-3}$} & \tentative{$8.86\times10^{-3}$} & \tentative{27} & \tentative{43\%}\\
mid resonance & \tentative{$2.25\times10^{-3}$} & \tentative{$6.65\times10^{-4}$} & \tentative{48} & \tentative{$2.65\times10^{-2}$} & \tentative{$5.15\times10^{-3}$} & \tentative{33} & \tentative{31\%}\\
weak slow & \tentative{$3.11\times10^{-4}$} & \tentative{$6.11\times10^{-4}$} & \tentative{46} & \tentative{$1.63\times10^{-4}$} & \tentative{$7.29\times10^{-4}$} & \tentative{31} & \tentative{33\%}\\
fast fast & \tentative{$3.99\times10^{-2}$} & \tentative{$7.84\times10^{-3}$} & \tentative{48} & \multicolumn{4}{c}{\tentative{pending}}\\
\end{tabular}
\end{ruledtabular}
\end{table}

The result is a cost--accuracy trade rather than a dominance relation.  This
route holds roughly a third fewer coordinates on every drive and is less
accurate in energy on five of six, by factors between three and twenty; each
method is reported at one setting of its own accuracy threshold, so these are
two points on a common trade-off rather than a statement about which frontier
is better.

\noindent\begin{minipage}{\columnwidth}
\begin{center}
\IfFileExists{figures/paper_ii_competitor_doublon_error_strongfast.pdf}{\includegraphics[width=0.98\columnwidth]{figures/paper_ii_competitor_doublon_error_strongfast.pdf}}{}
\end{center}
\captionof{figure}{Doublon-fraction deviation from the exact trajectory, strong fast drive.  \tentative{Preliminary.}}
\label{fig:overlay_strongfast}
\end{minipage}
\par\addvspace{1.0ex}

\noindent\begin{minipage}{\columnwidth}
\begin{center}
\IfFileExists{figures/paper_ii_competitor_doublon_error_fastweak.pdf}{\includegraphics[width=0.98\columnwidth]{figures/paper_ii_competitor_doublon_error_fastweak.pdf}}{}
\end{center}
\captionof{figure}{Doublon-fraction deviation from the exact trajectory, fast weak drive.  \tentative{Preliminary.}}
\label{fig:overlay_fastweak}
\end{minipage}
\par\addvspace{1.0ex}

\noindent\begin{minipage}{\columnwidth}
\begin{center}
\IfFileExists{figures/paper_ii_competitor_doublon_error_slowstrong.pdf}{\includegraphics[width=0.98\columnwidth]{figures/paper_ii_competitor_doublon_error_slowstrong.pdf}}{}
\end{center}
\captionof{figure}{Doublon-fraction deviation from the exact trajectory, slow strong drive.  \tentative{Preliminary.}}
\label{fig:overlay_slowstrong}
\end{minipage}
\par\addvspace{1.0ex}

\noindent\begin{minipage}{\columnwidth}
\begin{center}
\IfFileExists{figures/paper_ii_competitor_doublon_error_midres.pdf}{\includegraphics[width=0.98\columnwidth]{figures/paper_ii_competitor_doublon_error_midres.pdf}}{}
\end{center}
\captionof{figure}{Doublon-fraction deviation from the exact trajectory, mid resonance drive.  \tentative{Preliminary.}}
\label{fig:overlay_midres}
\end{minipage}
\par\addvspace{1.0ex}

\noindent\begin{minipage}{\columnwidth}
\begin{center}
\IfFileExists{figures/paper_ii_competitor_doublon_error_weakslow.pdf}{\includegraphics[width=0.98\columnwidth]{figures/paper_ii_competitor_doublon_error_weakslow.pdf}}{}
\end{center}
\captionof{figure}{Doublon-fraction deviation from the exact trajectory, weak slow drive.  \tentative{Preliminary.}}
\label{fig:overlay_weakslow}
\end{minipage}
\par\addvspace{1.0ex}

\noindent\begin{minipage}{\columnwidth}
\begin{center}
\IfFileExists{figures/paper_ii_competitor_doublon_error_fastfast.pdf}{\includegraphics[width=0.98\columnwidth]{figures/paper_ii_competitor_doublon_error_fastfast.pdf}}{}
\end{center}
\captionof{figure}{Doublon-fraction deviation from the exact trajectory, fast fast drive.  \tentative{Preliminary.}}
\label{fig:overlay_fastfast}
\end{minipage}
\par\addvspace{1.0ex}

\subsection{Integrator dependence of the shared numerical method}
\label{sec:integrator_ablation}

Because every arm in this comparison runs one inner numerical method, the choice
of that method is a property of the comparison rather than of any policy, and it
is reported here rather than assumed.  Table~\ref{tab:integrator_ablation}
varies only the integrator, holding the seed, drive, grid, Tikhonov
regularization ($\xi=10^{-6}$), candidate pool, and step control fixed.

\begin{table}[t]
\caption{Integrator dependence at fixed regularization and step control, adaptive-append policy on the driven Hubbard--Holstein seed to $t=10$.  The state-motion step guard is active and satisfied in every cell --- no step exhausted its subdivision budget --- so the difference is the integrator order alone.  \tentative{Two drives; the remaining four and the other policies are in progress.}}
\label{tab:integrator_ablation}
\begin{ruledtabular}
\begin{tabular}{llcc}
Drive & Integrator & $\overline{|\Delta E|}$ & $\overline{|\Delta n_{\rm d}|}$\\
\colrule
fast weak & fourth order & \tentative{$2.3\times10^{-4}$} & \tentative{$7.4\times10^{-4}$}\\
fast weak & Euler & \tentative{$2.6\times10^{-2}$} & \tentative{$6.3\times10^{-3}$}\\
strong fast & fourth order & \tentative{$3.8\times10^{-3}$} & \tentative{$1.0\times10^{-3}$}\\
strong fast & Euler & \tentative{$3.8\times10^{-2}$} & \tentative{$1.2\times10^{-2}$}\\
\end{tabular}
\end{ruledtabular}
\end{table}

Euler costs between one and two orders of magnitude in trajectory accuracy at
this reporting grid, and the step-size guard does not recover it: the guard was
satisfied at every checkpoint in all four cells, so the residual difference is
truncation error of the integrator itself rather than an uncontrolled step.
This bears on how the comparison should be read.  A first-order integrator is
what the adaptive-append source specifies, and it is cheaper per step, so a
comparison run entirely in Euler is faithful to that source and inexpensive; but
it places every policy in a regime where integrator error is comparable to or
larger than the structural differences under study, which is the regime in which
a structural claim is least legible.  We therefore report the integrator
alongside every result rather than fixing one, and we do not select it by which
choice favours the policy proposed here.

\subsection{Adaptive-append comparison}
\label{sec:comparator}

The closest adaptive-ansatz comparator is adaptive variational quantum dynamics simulation \cite{Yao2021}, which grows the ansatz during evolution using the McLachlan distance
\begin{equation}
L_k^2(J)=2\left[\|\bar b_k\|_2^2-Q_k(J)\right]_+
\label{eq:mclachlan_distance}
\end{equation}
as its trigger: while $L_k^2$ exceeds a fixed cut, it appends the pool operator whose zero-angle addition maximally reduces $L_k^2$, at the end of the circuit, and never deletes.  We implement that decision rule on the same checkpoint geometry, regularized inverse, solve repair, integrator, seed, drive, and time grid used by the exchange route, so the comparison isolates the structural rule rather than the numerical realization.  This is a comparison against the AVQDS \emph{rule}, not against its published implementation: the comparator inherits stabilization machinery that AVQDS as published does not specify, including the state-motion and temporal-kink subdivision guards of Sec.~\ref{sec:numerical_stabilization} and the regularized inverse.  We therefore make no claim to reproduce published AVQDS trajectories, and the growth-only ablation of our own route --- which retains positioned insertion at commutation-reduced cuts, cost-weighted ranking, and certification with refit, and differs from the full route only in that removal is disabled --- is reported separately from it rather than as the same object.  We adopt the comparator's append \emph{condition} for the exchange route's own insertions, since growth is not this work's contribution: insertions are considered when $L_k^2$ exceeds the cut and are repeated within one checkpoint until it is met or no candidate improves it.  The route's normalized residual-ratio gate is retired: under it the same route reached only $1.26\times10^{-1}$ mean energy error and $2.8\times10^{-2}$ on the doublon fraction at $67$ coordinates on this drive, against $2.1\times10^{-2}$ and $3.3\times10^{-3}$ at $37$ under the adopted condition.  The reason is measured rather than stylistic: a normalized gate is small precisely while the state is still accurate, so it defers growth past the window in which growth is cheap, and on this drive the comparator completed its growth by $t\approx0.9$ and held $46$ coordinates thereafter while the residual-gated route was still adding coordinates at $t\approx2$ with $L_k^2$ four orders of magnitude larger and a state already off the exact trajectory.  Under the adopted condition, the comparator considers only end-appended insertions, whereas the exchange route ranks complete generalized exchange patches by the same signed realized-drift change of Eq.~(\ref{eq:signed_patch_drift}); pure deletion and positioned insertion are its empty-component limits, and two nonempty components give true exchange.

Under this condition accuracy debt and cost relief are ordered rather than blended.  While $L_k^2$ exceeds the cut, signed realized-drift change is the primary score over the full generalized-exchange family; after the cut is met, the cost-aware utility of Eq.~(\ref{eq:operating_score}) resumes on its $B_k^{\oplus}=\varnothing$ face.  This distinction is necessary because the clipped deletion loss maps a helpful deletion and a harmless deletion to the same zero, after which the reciprocal cost-over-loss utility diverges as $1/\eps_L$.  Ranking by that composite utility during debt therefore spends the finalist budget on inexpensive deletions even when insertions would reduce the distance.  The insertion-only debt policy used for the provisional one-drive comparison below prevents that inversion by withdrawing removal, but it is retained as an ablation rather than the operating rule of Eqs.~(\ref{eq:signed_patch_drift})--(\ref{eq:debt_score}).

\begin{table}[t]
\caption{Adaptive-append comparison on one driven Hubbard--Holstein seed ($L=2$, $n_{\rm ph}^{\max}=1$) under the strong fast drive ($A=2.4$, $\omega=3$) to $t=10$ over $251$ checkpoints.  All policies share the seed, drive, reporting grid, inverse policy, solve repair, fourth-order integration, and the full deduplicated candidate pool of $125$ words; errors are against the exact time-ordered trajectory of the same truncated Hamiltonian.  The exchange row uses the insertion-only accuracy-debt ablation and is not evidence for the signed-drift operating rule of Eq.~(\ref{eq:debt_score}).  The pool is no longer capped below its own size, so the comparator reaches its McLachlan-distance cut by its own stopping rule rather than by a per-checkpoint budget, and no budget is imposed on it.  \tentative{One drive and one seed; the six-drive signed-drift validation is in progress.}}
\label{tab:comparator_hh_snake}
\begin{ruledtabular}
\begin{tabular}{lccc}
Policy & $\overline{|\Delta E|}$ & $\overline{|\Delta n_{\rm d}|}$ & $N_{\rm param}^{\rm final}$\\
\colrule
adaptive append \cite{Yao2021} & \tentative{$9.1\times10^{-3}$} & \tentative{$2.7\times10^{-3}$} & \tentative{46}\\
append-only (ablation) & \tentative{$5.3\times10^{-2}$} & \tentative{$9.8\times10^{-3}$} & \tentative{44}\\
exchange (sequenced debt) & \tentative{$2.1\times10^{-2}$} & \tentative{$3.3\times10^{-3}$} & \tentative{37}\\
\end{tabular}
\end{ruledtabular}
\end{table}

\noindent\begin{minipage}{\columnwidth}
\begin{center}
\IfFileExists{figures/paper_ii_obs_doublon_hh_strongfast_t10.pdf}{\includegraphics[width=0.98\columnwidth]{figures/paper_ii_obs_doublon_hh_strongfast_t10.pdf}}{}
\end{center}
\captionof{figure}{Doublon fraction to $t=10$ under the strong fast drive ($A=2.4$, $\omega=3$), Hubbard--Holstein seed, against the exact time-ordered trajectory.  Mean absolute deviations are \tentative{$3.3\times10^{-3}$} (insertion-only-debt exchange ablation, $37$ coordinates) and \tentative{$2.7\times10^{-3}$} (adaptive append, $46$).  \tentative{One drive; the signed-drift exchange validation is in progress.}}
\label{fig:obs_doublon_strongfast}
\end{minipage}
\par\addvspace{1.0ex}

\noindent\begin{minipage}{\columnwidth}
\begin{center}
\IfFileExists{figures/paper_ii_obs_occupations_hh_strongfast_t10.pdf}{\includegraphics[width=0.98\columnwidth]{figures/paper_ii_obs_occupations_hh_strongfast_t10.pdf}}{}
\end{center}
\captionof{figure}{Site occupations to $t=10$ under the same drive.  Both sites oscillate through the full horizon, so the occupations rather than $\langle H\rangle$ carry the dynamical content here; the exact curves are overlaid by both policies for the whole trajectory.  \tentative{One drive.}}
\label{fig:obs_occupations_strongfast}
\end{minipage}
\par\addvspace{1.0ex}

\noindent\begin{minipage}{\columnwidth}
\begin{center}
\IfFileExists{figures/paper_ii_obs_doublon_error_hh_strongfast_t10.pdf}{\includegraphics[width=0.98\columnwidth]{figures/paper_ii_obs_doublon_error_hh_strongfast_t10.pdf}}{}
\end{center}
\captionof{figure}{Deviation of the doublon fraction from the exact trajectory, same runs as Fig.~\ref{fig:obs_doublon_strongfast}.  The raw observable cannot separate the policies because all of them track it; the deviation can.  The append-only ablation sits near $10^{-2}$ for the whole horizon, while the insertion-only-debt exchange ablation and the AVQDS comparator run an order of magnitude below it and interleave, neither dominating.  \tentative{One drive.}}
\label{fig:obs_doublon_error_strongfast}
\end{minipage}
\par\addvspace{1.0ex}

\noindent\begin{minipage}{\columnwidth}
\begin{center}
\IfFileExists{figures/paper_ii_obs_energy_error_hh_strongfast_t10.pdf}{\includegraphics[width=0.98\columnwidth]{figures/paper_ii_obs_energy_error_hh_strongfast_t10.pdf}}{}
\end{center}
\captionof{figure}{Absolute energy error against the exact trajectory.  The discrete step in the insertion-only-debt exchange ablation near $t\approx8.6$ is the stabilization defect discussed in the text: it carries no structural edit, and it dominates that run's mean energy error while leaving the observables of Figs.~\ref{fig:obs_doublon_strongfast}--\ref{fig:obs_occupations_strongfast} essentially unchanged.  \tentative{One drive.}}
\label{fig:obs_energy_error_strongfast}
\end{minipage}
\par\addvspace{1.0ex}

\noindent\begin{minipage}{\columnwidth}
\begin{center}
\IfFileExists{figures/paper_ii_comparator_params_hh_strongfast.pdf}{\includegraphics[width=0.98\columnwidth]{figures/paper_ii_comparator_params_hh_strongfast.pdf}}{}
\end{center}
\captionof{figure}{Active support size under the three policies of Table~\ref{tab:comparator_hh_snake}.  The AVQDS comparator and the append-only ablation both grow monotonically, to $46$ and $44$ coordinates, because neither rule contains a removal mechanism; the insertion-only-debt exchange ablation front-loads growth under the same append condition and then contracts to $37$ after the distance cut is met, including a late contraction near $t\approx8.5$.  \tentative{One drive.}}
\label{fig:comparator_params}
\end{minipage}
\par\addvspace{1.0ex}

On the observables the two provisional policies are equivalent: the insertion-only-debt exchange ablation tracks the doublon fraction to $3.3\times10^{-3}$ against the comparator's $2.7\times10^{-3}$, and both overlay the exact site occupations through the full horizon (Figs.~\ref{fig:obs_doublon_strongfast}--\ref{fig:obs_occupations_strongfast}), which oscillate strongly to $t=10$ and are the discriminating observables on this drive.  The ablation holds $37$ coordinates against $46$, a $20\%$ smaller support, because it removes coordinates after the accuracy debt is paid while the comparator's rule contains no removal mechanism.  This comparison supports the bounded claim for that ablation only; the signed-drift operating route awaits the matched six-drive validation.

The append-only ablation, run under the same insertion condition and differing from the insertion-only-debt exchange ablation only in that removal is disabled, is worse on both axes: $9.8\times10^{-3}$ on the doublon fraction at $44$ coordinates, three times the exchange ablation's deviation on seven more coordinates.  Removal is therefore not merely a cost device on this drive; forbidding it costs accuracy as well as support, which is consistent with the conditioning relief that removal provides.

The energy is the one axis where the two separate, and the separation is not distributed over the trajectory: the insertion-only-debt exchange ablation holds $\sim\!1.5\times10^{-2}$ from $t\approx2$ to $t\approx8.6$ and then takes a single discrete step to $\sim\!8\times10^{-2}$, which dominates its mean.  That step carries no structural edit; the state-motion guard is active and escalating at it, to depth six and $64$ substeps with per-substep prospective motion $1.8\times10^{-3}$ against a $10^{-2}$ cap.  The guard bounds displacement per substep, so subdividing further shrinks each substep while permitting more of them and leaves the displacement accumulated across one checkpoint unbounded.  We report this as an open defect of the stabilization layer rather than evidence about the revised signed-drift structural rule, and note that it is nearly invisible in the observables the method targets: $\langle H\rangle$ resolves it because the truncated spectrum spans $[-0.4,\,10.4]$.  Three qualifications bound the evidence: one seed, a short horizon, and an insertion-only-debt policy superseded at the method level pending matched validation of Eq.~(\ref{eq:debt_score}).

\subsection{Factorial interactions}

The complete mechanism study crosses the structural policy arm (append-only, prune-only, exchange without refit, exchange with refit) with state-space numerical repair while retaining the same nominal inverse and reporting grid.  Insertion changes the dimension and retained spectrum presented to the solve; deletion can remove redundant or conditioning-toxic modes; repair determines how the selected support is realized over a finite step.  The crossed comparison therefore reports main effects and policy--repair interactions for trajectory error, integrated structural miss, unsupported checkpoints, final support, retained conditioning, measured-candidate work, and compiled resources.  A deliberately disabled numerical stabilizer is treated as an ablation, not as the canonical operating condition.

\subsection{Driven Hubbard--Holstein survey}

\structlabel{Main trajectory layout}
The primary results are organized as four one-column Hubbard--Holstein regime blocks.  Each block overlays the exact-diagonalization reference trajectory with AP-McLachlan dynamics and three pinned Qiskit-community baselines: TrotterQRTE, PVQD, and VarQRTE.  The plotted observable is the instantaneous total-energy expectation under the driven Hamiltonian; exact diagonalization is a diagnostic reference trajectory, not an input to the checkpoint policy.  The cost table below each trajectory uses the same Qiskit-compiled resource convention as the static-ansatz paper.  The four main regime blocks use the selected SNAKE plateau seed; append-only versus SNAKE seed effects are separated in Table~\ref{tab:dyn_seed_track_comparison}.

\noindent\begin{minipage}{\columnwidth}
\inlinetablecaption{tab:dyn_hh_weak_weak_cost}{Weak--weak Hubbard--Holstein trajectory and Qiskit-cost comparison.  The table is structurally fixed; red entries are diagnostic placeholders from a corrected energy-parity smoke.  The run uses candidate class settings and a three-sample grid, so these values are included to show structure rather than to support final claims.}
\begin{center}
\scriptsize
\setlength{\tabcolsep}{0.4pt}
\renewcommand{\arraystretch}{1.02}
\begin{ruledtabular}
\begin{tabular}{p{0.25\columnwidth}ccccccc}
Method & $\bar\epsilon_E$ & $\epsilon_{\obs}$ & $1-F$ & $N_{\twq}$ & $D_{\twq}$ & $D_c$ & $S$\\
\colrule
AP-McLachlan & \tentative{6.3e-6} & \tentative{1.2e-5} & \tentative{--} & \tentative{93} & \tentative{87} & \tentative{266} & \tentative{--}\\
TrotterQRTE & \tentative{1.5e-14} & \tentative{4.1e-14} & \tentative{2.8e-14} & \tentative{534} & \tentative{534} & \tentative{816} & \tentative{--}\\
PVQD & \tentative{6.8e-5} & \tentative{1.0e-2} & \tentative{7.3e-3} & \tentative{318} & \tentative{318} & \tentative{534} & \tentative{--}\\
VarQRTE & \tentative{3.8e-3} & \tentative{5.1e-3} & \tentative{1.5e-3} & \tentative{318} & \tentative{318} & \tentative{534} & \tentative{--}\\
\end{tabular}
\end{ruledtabular}
\end{center}
\end{minipage}

\par\addvspace{1.5ex}
\noindent\begin{minipage}{\columnwidth}
\inlinetablecaption{tab:dyn_hh_strong_weak_cost}{Strong--weak Hubbard--Holstein trajectory and Qiskit-cost comparison.}
\begin{center}
\scriptsize
\setlength{\tabcolsep}{1.0pt}
\renewcommand{\arraystretch}{1.02}
\begin{ruledtabular}
\begin{tabular}{p{0.30\columnwidth}ccccccc}
Method & \metrichead{$\bar\epsilon_E$} & \metrichead{$\epsilon_{\obs}$} & \metrichead{$1-F$} & \metrichead{$N_{\twq}$} & \metrichead{$D_{\twq}$} & \metrichead{$D_c$} & \metrichead{$S$}\\
\colrule
AP-McLachlan & \tentative{--} & \tentative{--} & \tentative{--} & \tentative{--} & \tentative{--} & \tentative{--} & \tentative{--}\\
Qiskit TrotterQRTE & \tentative{--} & \tentative{--} & \tentative{--} & \tentative{--} & \tentative{--} & \tentative{--} & \tentative{--}\\
Qiskit PVQD & \tentative{--} & \tentative{--} & \tentative{--} & \tentative{--} & \tentative{--} & \tentative{--} & \tentative{--}\\
Qiskit VarQRTE & \tentative{--} & \tentative{--} & \tentative{--} & \tentative{--} & \tentative{--} & \tentative{--} & \tentative{--}\\
\end{tabular}
\end{ruledtabular}
\end{center}
\end{minipage}

\par\addvspace{1.5ex}
\noindent\begin{minipage}{\columnwidth}
\inlinetablecaption{tab:dyn_hh_weak_strong_cost}{Weak--strong Hubbard--Holstein trajectory and Qiskit-cost comparison.}
\begin{center}
\scriptsize
\setlength{\tabcolsep}{1.0pt}
\renewcommand{\arraystretch}{1.02}
\begin{ruledtabular}
\begin{tabular}{p{0.30\columnwidth}ccccccc}
Method & \metrichead{$\bar\epsilon_E$} & \metrichead{$\epsilon_{\obs}$} & \metrichead{$1-F$} & \metrichead{$N_{\twq}$} & \metrichead{$D_{\twq}$} & \metrichead{$D_c$} & \metrichead{$S$}\\
\colrule
AP-McLachlan & \tentative{--} & \tentative{--} & \tentative{--} & \tentative{--} & \tentative{--} & \tentative{--} & \tentative{--}\\
Qiskit TrotterQRTE & \tentative{--} & \tentative{--} & \tentative{--} & \tentative{--} & \tentative{--} & \tentative{--} & \tentative{--}\\
Qiskit PVQD & \tentative{--} & \tentative{--} & \tentative{--} & \tentative{--} & \tentative{--} & \tentative{--} & \tentative{--}\\
Qiskit VarQRTE & \tentative{--} & \tentative{--} & \tentative{--} & \tentative{--} & \tentative{--} & \tentative{--} & \tentative{--}\\
\end{tabular}
\end{ruledtabular}
\end{center}
\end{minipage}

\par\addvspace{1.5ex}
\noindent\begin{minipage}{\columnwidth}
\inlinetablecaption{tab:dyn_hh_strong_strong_cost}{Strong--strong Hubbard--Holstein trajectory and Qiskit-cost comparison.}
\begin{center}
\scriptsize
\setlength{\tabcolsep}{1.0pt}
\renewcommand{\arraystretch}{1.02}
\begin{ruledtabular}
\begin{tabular}{p{0.30\columnwidth}ccccccc}
Method & \metrichead{$\bar\epsilon_E$} & \metrichead{$\epsilon_{\obs}$} & \metrichead{$1-F$} & \metrichead{$N_{\twq}$} & \metrichead{$D_{\twq}$} & \metrichead{$D_c$} & \metrichead{$S$}\\
\colrule
AP-McLachlan & \tentative{--} & \tentative{--} & \tentative{--} & \tentative{--} & \tentative{--} & \tentative{--} & \tentative{--}\\
Qiskit TrotterQRTE & \tentative{--} & \tentative{--} & \tentative{--} & \tentative{--} & \tentative{--} & \tentative{--} & \tentative{--}\\
Qiskit PVQD & \tentative{--} & \tentative{--} & \tentative{--} & \tentative{--} & \tentative{--} & \tentative{--} & \tentative{--}\\
Qiskit VarQRTE & \tentative{--} & \tentative{--} & \tentative{--} & \tentative{--} & \tentative{--} & \tentative{--} & \tentative{--}\\
\end{tabular}
\end{ruledtabular}
\end{center}
\end{minipage}

\subsection{Append-only versus SNAKE seed tracks}

\noindent\begin{minipage}{\columnwidth}
\inlinetablecaption{tab:dyn_seed_track_comparison}{Seed-track comparison held separate from the trajectory overlays.  Each pair is selected at matched same-cutoff static energy error, then propagated with the same Hamiltonian, drive, time grid, observable set, Qiskit backend, and AP-McLachlan dynamics rule.  The selected depth may be a plateau prefix rather than the terminal static-ansatz depth.}
\begin{center}
\scriptsize
\setlength{\tabcolsep}{0.9pt}
\renewcommand{\arraystretch}{1.02}
\begin{ruledtabular}
\begin{tabular}{p{0.20\columnwidth}p{0.14\columnwidth}ccccc}
Regime & Seed & $|\Delta E_{\rm seed}|$ & $k_{\rm seed}$ & $N_{\twq}^{\rm prep}$ & $D_{\twq}^{\rm prep}$ & $\bar\epsilon_E^{\rm AP}$\\
\colrule
weak--weak & append & \tentative{--} & \tentative{--} & \tentative{--} & \tentative{--} & \tentative{--}\\
weak--weak & SNAKE & \tentative{--} & \tentative{--} & \tentative{--} & \tentative{--} & \tentative{--}\\
strong--weak & append & \tentative{--} & \tentative{--} & \tentative{--} & \tentative{--} & \tentative{--}\\
strong--weak & SNAKE & \tentative{--} & \tentative{--} & \tentative{--} & \tentative{--} & \tentative{--}\\
weak--strong & append & \tentative{--} & \tentative{--} & \tentative{--} & \tentative{--} & \tentative{--}\\
weak--strong & SNAKE & \tentative{--} & \tentative{--} & \tentative{--} & \tentative{--} & \tentative{--}\\
strong--strong & append & \tentative{--} & \tentative{--} & \tentative{--} & \tentative{--} & \tentative{--}\\
strong--strong & SNAKE & \tentative{--} & \tentative{--} & \tentative{--} & \tentative{--} & \tentative{--}\\
\end{tabular}
\end{ruledtabular}
\end{center}
\end{minipage}

\subsection{Legacy controller diagnostics (archived)}

Provisional diagnostic tables produced by the retired checkpoint-controller implementation have been moved to the repository archive (\texttt{MATH/paper\_details/archive/}) pending regeneration under the exchange selector on the same Hamiltonian and seed surfaces.  No quantity from those tables is cited by this manuscript.

\section{Discussion}
\label{sec:discussion}

The two-failure decomposition assigns distinct authority to distinct diagnostics. Structural miss measures tangent motion absent from the resolved manifold. Numerical miss measures supported motion lost by the realized inverse. The effective condition number measures susceptibility of the retained solve but does not determine the direction of repair. Parameter magnitude determines neither current deletion loss nor trajectory continuity, and exact trajectory error is unavailable to a measurement-compatible online policy.

The patch-then-realize organization also fixes the contribution boundary. Adaptive growth, zero-angle append, regularization, eigencutoffs, pseudoinverses, and local subdivision are established ingredients. The structural contribution is one generalized-exchange finalist family with removal and addition components. Stay, child-resolved cost-weighted append, and continuity-certified prune are its empty-component limits; two nonempty components give true exchange. The accepted patch is confirmed on the complete post-patch tangent support. The numerical contribution then acts on that fixed support through bidirectional inverse candidates, separation of inverse defects from excessive finite-step motion, least-intervention selection, and hysteretic release.

Prune is not merely post hoc compression. A redundant tangent direction can increase circuit depth, estimator burden, and retained-spectrum fragility. Removing it can improve numerical realization while the same-checkpoint state is intentionally preserved. Conversely, hardware expense or poor conditioning cannot authorize deletion when current grouped loss or continuity is large. This asymmetry is deliberate: a zero-angle append preserves the state automatically but must justify its vector-field and resource effect, whereas a deletion promises resource or conditioning relief but bears the stronger continuity burden.

The present study is limited to small, classically verifiable gate-based instances. Combinatorial candidate frontiers require bounded search; finite-shot uncertainty can perturb retained rank, Schur novelty, and deletion loss; and same-checkpoint continuity does not prove global no-harm. Final compiled ansatz cost, online estimator work, and full-horizon repeated-execution cost are therefore reported as separate resource axes. True exchange remains an empirical result only when a nonempty delete--append patch passes the audited joint commit record.

\subsection{Evidence scope and planned extensions}

The structural evidence reported here is bounded in four ways, each addressed by a planned production batch rather than by reinterpretation of the present runs.  Trajectories terminate at $t=1$, so accumulated realization error and support staleness are observed only in their onset; production horizons extend to $t\gtrsim5$ where both mechanisms and the comparator's own threshold reach their asymptotic regimes.  Each structural comparison uses one static-ansatz seed per redundancy class, so seed-to-seed variation is not yet quantified.  The Hamiltonian family is Hubbard--Holstein at $L=2$ with binary-aligned phonon cutoffs, and the driven survey spans the four coupling regimes of Sec.~\ref{sec:results} rather than the wider benchmark family used for the fixed-support comparators.  Reported observables are the total energy, the doublon fraction, and site occupations; correlation functions and phonon-sector observables are natural additions once the horizon is extended.

\section{Conclusion}
\label{sec:conclusion}

AP-McLachlan treats real-time variational simulation as the coordinated maintenance and numerical realization of a live tangent manifold. Its structural action is one jointly evaluated generalized exchange patch. Empty removal or addition components give stay, pure append, and pure prune as limiting cases, while two nonempty components give true exchange. The accepted support is then propagated through state-space solve repair and local subdivision when the realized McLachlan motion is numerically defective or too large for the proposed step. Exact reference trajectories remain excluded from both decisions. Quantitative regime-wide conclusions remain tied to the source-locked mechanism matrix and driven Hubbard--Holstein survey.
\begin{acknowledgments}
\placeholder{Add funding, compute resources, discussions, and advisor acknowledgments.}
\end{acknowledgments}

\appendix

\section{Technical realization details}
\label{app:technical_realization}

The main text states the physical diagnostics and accepted structural actions.  This appendix records the numerical coordinate conventions and controller details needed to reproduce their implementation.

\subsection{Retained support, damped inversion, and whitening}

Let the measured Gram matrix have retained eigendecomposition
\begin{equation}
G_kV_k=V_k\Lambda_k,
\qquad
\lambda_a>\tau_{\rm pinv}\lambda_{\max}.
\label{eq:app_retained_support}
\end{equation}
The exact supported inverse and the damped realization are
\begin{equation}
G_{k,\tau}^{\dagger}
=V_k\Lambda_k^{-1}V_k^\top,
\qquad
G_{k,\tau,\eps}^{\oplus}
=V_k(\Lambda_k+\eps_{\rm solve}I)^{-1}V_k^\top.
\label{eq:app_supported_inverse}
\end{equation}
The retained-spectrum diagnostic is
\begin{equation}
\kappa_{\rm eff}(G_k)
=\frac{\lambda_{\max}}{\lambda_{\min}^{\rm kept}}.
\label{eq:app_kappa_eff}
\end{equation}
As in the accepted-refit notation of Ref.~\cite{Strobel2026StaticAdapt}, define
\begin{equation}
W_k=V_k\Lambda_k^{-1/2},
\qquad
W_k^\top G_kW_k=I.
\label{eq:app_whitening}
\end{equation}
Thus \(W_k\) maps orthonormal retained tangent coordinates back to circuit-coordinate displacements.  For an incoming tangent block with current--candidate cross Gram \(B_k\) and candidate Gram \(C_k\), whitening gives
\begin{equation}
\widetilde B_k=W_k^\top B_k,
\qquad
\mathsf S_k=C_k-\widetilde B_k^\top\widetilde B_k,
\label{eq:app_whitened_schur}
\end{equation}
which is the retained-support novelty block of the incoming coordinates against the active span.  Whitening and damping stabilize the computation; the undamped \(G\)-based gain and loss in the main text remain the structural definitions.

\subsection{Numerical-repair candidate construction}

A numerical candidate specifies an inverse convention and, when needed, a local subdivision factor.  Numerical miss opens candidate inverses on both sides of the base regularization; excessive state displacement or temporal kink opens successively finer local subdivisions.  The base convention remains in the candidate set.  Hard-invalid candidates are removed, and the remaining candidates are tested against
\begin{equation}
\rho_{{\rm num},k}\le\tau_{\rm num},
\qquad
\delta_{{\rm step},k}\le\tau_\Delta,
\qquad
\eta_{{\rm time},k}^{\psi}\le\tau_{\rm kink}.
\label{eq:app_repair_acceptance}
\end{equation}
Among candidates satisfying Eq.~(\ref{eq:app_repair_acceptance}), the implementation minimizes a declared least-intervention score combining departure from the base convention, normalized defect sizes, and the soft retained-spectrum risk.  The condition number can penalize a fragile candidate but is a hard rejection only in an explicitly strict validation mode.  If no finite candidate passes, diagnostic runs may continue with the least-bad finite candidate and mark the checkpoint unsupported.

Same-checkpoint inverse candidates reuse \(G_k\), \(f_k\), and \(\|\bar b_k\|_2^2\).  In particular,
\begin{equation}
\begin{aligned}
\|\bar T_kv\|_2^2
&=v^\top G_kv,\\
\|\bar T_kv-\bar b_k\|_2^2
&=v^\top G_kv-2f_k^\top v+\|\bar b_k\|_2^2.
\end{aligned}
\label{eq:app_classical_quadratic_forms}
\end{equation}
so testing inverse candidates and the release metric in Eq.~(\ref{eq:repair_release}) is classical after the checkpoint data have been acquired.  Local subdivision creates new states and therefore requires new checkpoint evaluations.

\subsection{Diagnostic-only machinery (not part of the operating method)}
\label{app:diagnostic_only}

Everything in this subsection is retained for diagnosis of pathological
trajectories.  None of it is part of the method used for any result in this
paper: the operating numerical rule is the single
state-displacement subdivision of Sec.~\ref{sec:numerical_solve_repair}, and
the operating structural rule is Eqs.~(\ref{eq:operating_score}) and
(\ref{eq:debt_score}).  The material
is recorded because it is the evidence that the reduced configuration loses
nothing, not because a reader must apply it.



The repair lane opens only when the base realization is invalid or one of these measured defects is excessive:
\begin{equation}
g_{{\rm invalid},k}
\ \vee\
\rho_{{\rm num},k}>\tau_{\rm num}
\ \vee\
\delta_{{\rm step},k}>\tau_\Delta
\ \vee\
\eta_{{\rm time},k}^{\psi}>\tau_{\rm kink}.
\label{eq:repair_entry}
\end{equation}
Hard invalidity means nonfinite or dimensionally inconsistent checkpoint data, not merely a large diagnostic value.  The retained-spectrum condition number \(\kappa_{\rm eff}\) is intentionally absent from Eq.~(\ref{eq:repair_entry}).  It measures sensitivity of a candidate inverse and therefore modifies candidate risk; by itself it neither proves that motion was lost nor determines whether regularization should increase or decrease.

After a repaired inverse has been held, return to the base inverse is tested at the same checkpoint.  If \(v_k^{\rm rep}\) and \(v_k^{0}\) are the repaired and base velocities, respectively, then
\begin{equation}
\eta_{{\rm rel},k}^{\psi}
=
\frac{
(v_k^{\rm rep}-v_k^{0})^\top
G_k
(v_k^{\rm rep}-v_k^{0})
}{
\|\bar b_k\|_2^2+\eps_{\rm norm}
}.
\label{eq:repair_release}
\end{equation}
Equation~(\ref{eq:repair_release}) is computed classically from the already measured \(G_k\); it does not require a second quantum measurement of either velocity.  Release uses a stricter threshold and severity-dependent persistence so that a repair is not switched off immediately after a large accepted correction.

\paragraph*{Operating configuration.}  The candidate set above is the diagnostic superset.  Across every trajectory reported here --- 885 integration steps --- the inverse-candidate search left the base regularization in place on all but two steps, damping never moved from its base value, and of $1.15\times10^4$ evaluated candidates six were accepted; the conditioning, numerical-miss, and temporal-kink triggers opened searches that did not change the accepted solve.  Local subdivision is the repair mechanism that repeatedly altered a step, through both of its triggers: $372$ subdivisions opened on excessive prospective state displacement and $254$ on the temporal-kink diagnostic.  The reported runs therefore use a reduced configuration containing both subdivision triggers and the base inverse convention alone.  Substituting the reduced configuration for the full candidate search on the driven seed of Table~\ref{tab:comparator_hh_snake} changes the mean energy error from $3.4735\times10^{-3}$ to $3.4738\times10^{-3}$ at identical final support, while removing a quarter of the candidate solves.  The wider set is retained for diagnosis of pathological checkpoints, not as the operating point.

\subsection{Structural search and deletion history}

Candidate tangent overlaps are measured once per checkpoint and reused across the whole structural family: the positioned candidate tangents are assembled in one batched pass, the deletion-independent Gram and force blocks are cached, and every candidate solve is a submatrix solve of that cached geometry with memoized results.  Deletion branches reuse the base-support solve metadata, so the conditioning terms below require no additional eigendecompositions.

Outside accuracy debt, prune nomination retains Eq.~(\ref{eq:prune_loss}) as its damage authority.  Its implementation-level utility may include historical loss, conditioning relief, Schur degeneracy, and conditioning damage:
\begin{equation}
\begin{aligned}
S_k^{\ominus,\kappa}(B)
&=\frac{K_{t_k}(B)^{\alpha_{\rm prune}}}{
\mathcal L_k^{\ominus}(B)
+\lambda_{\rm hist}\bar{\mathcal L}_k^{\ominus,\rm hist}(B)
+\lambda_{\kappa}^{\rm dam}d_{\kappa}^{\rm dam}
+\eps_L}\\
&\quad\times
\Bigl(1+\lambda_{\kappa}^{\rm rel}d_{\kappa}^{\rm rel}
+\lambda_Sd_S\\[-2pt]
&\hspace{5.7em}
+\lambda_{\kappa}^{\rm hist}d_{\kappa,S}^{\rm hist}\Bigr).
\end{aligned}
\label{eq:app_prune_nomination}
\end{equation}
The conditioning and history terms in Eq.~(\ref{eq:app_prune_nomination}) are extensions carried at zero weight in every provisional run reported here, and no reported result depends on a nonzero $\lambda$.  During accuracy debt, Eq.~(\ref{eq:app_prune_nomination}) contributes only through the bounded secondary term of Eq.~(\ref{eq:debt_score}); the signed realized-drift change is primary.  After the distance cut is met, the cost-over-loss utility resumes.  Current patched solve geometry, reduced-state refit, ray continuity, velocity continuity, persistence, cooldown, and rollback retain commit authority.  Atom-level persistence survives append-only support changes for atoms that remain active; support-conditional loss and smoothness histories are invalidated when the support changes.

True exchange is enumerated directly: within each admitted deletion rung, every deletion branch is paired with every positioned insertion from its deletion-conditioned pool, with the insertion cuts remapped onto the post-deletion survivor word.  Every candidate is evaluated on
\[
J_k^{\mathcal B}
=(J_k\setminus I_{B,k}^{\ominus})\cup I_{B,k}^{\oplus},
\]
using conditional insertion gain against the post-delete support, conditional deletion loss against the patched support, the full patched solve, and Eqs.~(\ref{eq:patch_continuity})--(\ref{eq:patch_velocity_continuity}) at certification.  This joint evaluation is what distinguishes exchange from sequential append followed by prune.  Certification order and the per-level and per-branch admission bounds are as in the main text; the historical-loss term \(\bar{\mathcal L}_k^{\ominus,\rm hist}\) in Eq.~(\ref{eq:app_prune_nomination}) is a windowed mean of previously attempted deletion losses recorded per surviving coordinate, and the conditioning relief and damage terms \(d_{\kappa}^{\rm rel},d_{\kappa}^{\rm dam}\) are the positive and negative parts of the base-versus-branch log-condition shift of the retained solve.


\begin{thebibliography}{99}

\bibitem{Strobel2026StaticAdapt}
J. Strobel, J. Simoni, and Y. Ping,
``Geometric and Resource-Aware Adaptive Ansatz Construction for Fermion--Boson Hamiltonian Systems,''
manuscript in preparation (2026).

\bibitem{Weisse2008}
A. Wei{\ss}e and H. Fehske,
``Exact diagonalization techniques,''
in \emph{Computational Many-Particle Physics},
edited by H. Fehske, R. Schneider, and A. Wei{\ss}e,
Lecture Notes in Physics Vol. 739 (Springer, Berlin, Heidelberg, 2008), pp. 529--544.

\bibitem{Troyer2005}
M. Troyer and U.-J. Wiese,
``Computational complexity and fundamental limitations to fermionic quantum Monte Carlo simulations,''
\emph{Physical Review Letters} \textbf{94}, 170201 (2005).

\bibitem{Schollwock2011}
U. Schollw\"ock,
``The density-matrix renormalization group in the age of matrix product states,''
\emph{Annals of Physics} \textbf{326}, 96 (2011).

\bibitem{NocedalWright2006}
J. Nocedal and S.~J. Wright,
\emph{Numerical Optimization}, 2nd ed.
(Springer, New York, 2006).

\bibitem{IglewiczHoaglin1993}
B. Iglewicz and D.~C. Hoaglin,
\emph{How to Detect and Handle Outliers},
ASQC Basic References in Quality Control: Statistical Techniques, Vol.~16
(ASQC Quality Press, Milwaukee, WI, 1993).

\bibitem{Burbidge1988IHS}
J.~B. Burbidge, L. Magee, and A.~L. Robb,
``Alternative transformations to handle extreme values of the dependent variable,''
\emph{Journal of the American Statistical Association} \textbf{83}, 123--127 (1988).

\bibitem{Ikhtiarudin2025ShotEfficient}
A. Ikhtiarudin, G.~K. Sunnardianto, F. Fathurrahman, M.~K. Agusta, and H. Dipojono,
``Shot-Efficient ADAPT-VQE via Reused Pauli Measurements and Variance-Based Shot Allocation,''
\emph{arXiv:2507.16879} (2025).

\bibitem{Verteletskyi2020CliqueCover}
V. Verteletskyi, T.-C. Yen, and A.~F. Izmaylov,
``Measurement optimization in the variational quantum eigensolver using a minimum clique cover,''
\emph{Journal of Chemical Physics} \textbf{152}, 124114 (2020).

\bibitem{Yen2020CompatibleMeasurements}
T.-C. Yen, V. Verteletskyi, and A.~F. Izmaylov,
``Measuring all compatible operators in one series of single-qubit measurements using unitary transformations,''
\emph{Journal of Chemical Theory and Computation} \textbf{16}, 2400--2409 (2020).

\bibitem{YuanLin2006}
M. Yuan and Y. Lin,
``Model selection and estimation in regression with grouped variables,''
\emph{Journal of the Royal Statistical Society: Series B} \textbf{68}, 49 (2006).

\bibitem{FrischWaugh1933}
R. Frisch and F.~V. Waugh,
``Partial time regressions as compared with individual trends,''
\emph{Econometrica} \textbf{1}, 387 (1933).

\bibitem{Lovell1963}
M.~C. Lovell,
``Seasonal adjustment of economic time series and multiple regression analysis,''
\emph{Journal of the American Statistical Association} \textbf{58}, 993 (1963).

\bibitem{BertsimasKingMazumder2016}
D. Bertsimas, A. King, and R. Mazumder,
``Best subset selection via a modern optimization lens,''
\emph{The Annals of Statistics} \textbf{44}, 813 (2016).

\bibitem{Natarajan1995SparseApprox}
B.~K. Natarajan,
``Sparse approximate solutions to linear systems,''
\emph{SIAM Journal on Computing} \textbf{24}, 227 (1995).

\bibitem{DasKempe2011}
A. Das and D. Kempe,
``Submodular meets spectral: Greedy algorithms for subset selection, sparse approximation and dictionary selection,''
in \emph{Proceedings of the 28th International Conference on Machine Learning} (2011).

\bibitem{Elenberg2018WeakSubmodularity}
E.~R. Elenberg, R. Khanna, A.~G. Dimakis, and S. Negahban,
``Restricted strong convexity implies weak submodularity,''
\emph{The Annals of Statistics} \textbf{46}, 3539 (2018).

\bibitem{Efron1979}
B. Efron,
``Bootstrap methods: Another look at the jackknife,''
\emph{The Annals of Statistics} \textbf{7}, 1 (1979).

\bibitem{MeinshausenBuhlmann2010}
N. Meinshausen and P. B\"uhlmann,
``Stability selection,''
\emph{Journal of the Royal Statistical Society: Series B} \textbf{72}, 417 (2010).

\bibitem{Grimsley2019}
H.~R. Grimsley, S.~E. Economou, E. Barnes, and N.~J. Mayhall,
``An adaptive variational algorithm for exact molecular simulations on a quantum computer,''
\emph{Nature Communications} \textbf{10}, 3007 (2019).

\bibitem{Tang2021}
H.~L. Tang, V.~O. Shkolnikov, G.~S. Barron, H.~R. Grimsley, N.~J. Mayhall, E. Barnes, and S.~E. Economou,
``Qubit-ADAPT-VQE: An adaptive algorithm for constructing hardware-efficient ans\"atze on a quantum processor,''
\emph{PRX Quantum} \textbf{2}, 020310 (2021).

\bibitem{Yordanov2021}
Y.~S. Yordanov, V. Armaos, C.~H.~W. Barnes, and D.~R.~M. Arvidsson-Shukur,
``Qubit-excitation-based adaptive variational quantum eigensolver,''
\emph{Communications Physics} \textbf{4}, 228 (2021).

\bibitem{Feniou2023}
C. Feniou, M. Hassan, D. Traor\'e, E. Giner, Y. Maday, and J.-P. Piquemal,
``Overlap-ADAPT-VQE: Practical quantum chemistry on quantum computers via overlap-guided compact ans\"atze,''
\emph{Communications Physics} \textbf{6}, 192 (2023).

\bibitem{Anastasiou2024}
P.~G. Anastasiou, Y. Chen, N.~J. Mayhall, E. Barnes, and S.~E. Economou,
``TETRIS-ADAPT-VQE: An adaptive algorithm that yields shallower, denser circuit ans\"atze,''
\emph{Physical Review Research} \textbf{6}, 013254 (2024).

\bibitem{Ramoa2025}
M. Ram\^oa, P.~G. Anastasiou, L.~P. Santos, N.~J. Mayhall, E. Barnes, and S.~E. Economou,
``Reducing the resources required by ADAPT-VQE using coupled exchange operators and improved subroutines,''
\emph{npj Quantum Information} \textbf{11}, 86 (2025).

\bibitem{McLachlan1964}
A.~D. McLachlan,
``A variational solution of the time-dependent Schr\"odinger equation,''
\emph{Molecular Physics} \textbf{8}, 39 (1964).

\bibitem{LiBenjamin2017VQS}
Y. Li and S.~C. Benjamin,
``Efficient variational quantum simulator incorporating active error minimization,''
\emph{Physical Review X} \textbf{7}, 021050 (2017).

\bibitem{Yuan2019VQS}
X. Yuan, S. Endo, Q. Zhao, Y. Li, and S.~C. Benjamin,
``Theory of variational quantum simulation,''
\emph{Quantum} \textbf{3}, 191 (2019).

\bibitem{Endo2020GeneralProcesses}
S. Endo, J. Sun, Y. Li, S.~C. Benjamin, and X. Yuan,
``Variational quantum simulation of general processes,''
\emph{Physical Review Letters} \textbf{125}, 010501 (2020).

\bibitem{Benedetti2021TimeEvolution}
M. Benedetti, M. Fiorentini, and M. Lubasch,
``Hardware-efficient variational quantum algorithms for time evolution,''
\emph{Physical Review Research} \textbf{3}, 033083 (2021).

\bibitem{Miessen2021SpinBoson}
A. Miessen, P.~J. Ollitrault, and I. Tavernelli,
``Quantum algorithms for quantum dynamics: A performance study on the spin-boson model,''
\emph{Physical Review Research} \textbf{3}, 043212 (2021).

\bibitem{Barison2021}
S. Barison, F. Vicentini, and G. Carleo,
``An efficient quantum algorithm for the time evolution of parameterized circuits,''
\emph{Quantum} \textbf{5}, 512 (2021).

\bibitem{Yao2021}
Y.-X. Yao, N. Gomes, F. Zhang, C.-Z. Wang, K.-M. Ho, T. Iadecola, and P. P. Orth,
``Adaptive variational quantum dynamics simulations,''
\emph{PRX Quantum} \textbf{2}, 030307 (2021).

\bibitem{Linteau2024}
D. Linteau, S. Barison, N.~H. Lindner, and G. Carleo,
``Adaptive projected variational quantum dynamics,''
\emph{Physical Review Research} \textbf{6}, 023130 (2024).

\bibitem{Zhang2025}
F. Zhang, N. Gomes, Y.-X. Yao, P. P. Orth, C.-Z. Wang, K.-M. Ho, and T. Iadecola,
``Adaptive variational quantum dynamics simulations with compressed circuits,''
\emph{Physical Review B} \textbf{111}, 094310 (2025).

\bibitem{Suzuki1985}
M. Suzuki,
``Decomposition formulas of exponential operators and Lie exponentials with some applications to quantum mechanics and statistical physics,''
\emph{Journal of Mathematical Physics} \textbf{26}, 601 (1985).

\bibitem{Suzuki1990}
M. Suzuki,
``Fractal decomposition of exponential operators with applications to many-body theories and Monte Carlo simulations,''
\emph{Physics Letters A} \textbf{146}, 319 (1990).

\bibitem{ChildsSu2019ProductFormula}
A.~M. Childs and Y. Su,
``Nearly optimal lattice simulation by product formulas,''
\emph{Physical Review Letters} \textbf{123}, 050503 (2019).

\bibitem{ChildsOstranderSu2019RandomizedPF}
A.~M. Childs, A. Ostrander, and Y. Su,
``Faster quantum simulation by randomization,''
\emph{Quantum} \textbf{3}, 182 (2019).

\bibitem{Campbell2019QDRIFT}
E. Campbell,
``Random compiler for fast Hamiltonian simulation,''
\emph{Physical Review Letters} \textbf{123}, 070503 (2019).

\bibitem{Macridin2018}
A. Macridin, P. Spentzouris, J. Amundson, and R. Harnik,
``Electron-phonon systems on a universal quantum computer,''
\emph{Physical Review Letters} \textbf{121}, 110504 (2018).

\bibitem{Macridin2018PRA}
A. Macridin, P. Spentzouris, J. Amundson, and R. Harnik,
``Digital quantum computation of fermion-boson interacting systems,''
\emph{Physical Review A} \textbf{98}, 042312 (2018).

\bibitem{Sawaya2020}
N.~P.~D. Sawaya, T. Menke, T.~H. Kyaw, S. Johri, A. Aspuru-Guzik, and G.~G. Guerreschi,
``Resource-efficient digital quantum simulation of $d$-level systems for photonic, vibrational, and spin-$s$ Hamiltonians,''
\emph{npj Quantum Information} \textbf{6}, 49 (2020).

\bibitem{Li2023}
W. Li, J. Ren, S. Huai, T. Cai, Z. Shuai, and S. Zhang,
``Efficient quantum simulation of electron-phonon systems by variational basis state encoder,''
\emph{Physical Review Research} \textbf{5}, 023046 (2023).

\bibitem{Denner2023}
M.~M. Denner, A. Miessen, H. Yan, I. Tavernelli, T. Neupert, E. Demler, and Y. Wang,
``A hybrid quantum-classical method for electron-phonon systems,''
\emph{Communications Physics} \textbf{6}, 233 (2023).

\end{thebibliography}

\end{document}
